\chapter{Gastroenterology}

\medterm{Peritonitis}-- Inflammation or infection of the peritoneum, commonly caused by a perforated abdominal organ, appendicitis, abdominal trauma, or infection of ascitic fluid. Severe abdominal pain, guarding, fever, vomiting, or shock requires urgent medical evaluation.

\medterm{Nausea}-- An unpleasant sensation of an impending need to vomit, often accompanied by autonomic symptoms such as salivation, sweating, pallor, or dizziness. Common causes include gastrointestinal illness, pregnancy, medication effects, migraine, vestibular disorders, metabolic disease, and serious abdominal or neurologic conditions.

\medterm{Rectal Bleeding}-- Passage of bright-red or dark blood from the rectum, caused by hemorrhoids, fissures, polyps, inflammatory bowel disease, diverticular disease, or cancer. Persistent, heavy, or unexplained bleeding requires evaluation.

\medterm{Abdominal Bloating}-- A sensation of fullness or visible enlargement of the abdomen, commonly caused by intestinal gas, constipation, food intolerance, or functional gastrointestinal disorders. Persistent or severe bloating may require evaluation for organic disease.

\medterm{Belching}-- The release of swallowed air or gas from the stomach through the mouth. Frequent belching may occur with aerophagia, gastroesophageal reflux, functional dyspepsia, or carbonated-drink consumption.

\medterm{Anorectal Abscess}-- A localized collection of pus in the tissues around the anus or rectum, usually caused by infection of an anal gland. It commonly causes severe anal pain, swelling, fever, and difficulty sitting. Prompt incision and drainage is the standard treatment.

\medterm{Anoscopy}-- Examination of the anal canal with a short lighted instrument called an anoscope. It is used to evaluate rectal bleeding, hemorrhoids, fissures, fistulas, and other anorectal disorders.

\medterm{Peptic Ulcer}-- An open sore in the stomach or duodenal lining caused most often by Helicobacter pylori infection or nonsteroidal anti-inflammatory drug use. It may cause burning upper-abdominal pain, bleeding, perforation, or obstruction. Also called peptic ulcer disease.

\medterm{Bacterial Gastroenteritis}-- Acute inflammation of the stomach and intestines caused by pathogenic bacteria such as Salmonella, Campylobacter, Shigella, or pathogenic Escherichia coli. It commonly causes diarrhea, abdominal cramps, fever, and sometimes bloody stools or vomiting.

\medterm{Indigestion}-- Dyspepsia characterized by recurrent upper-abdominal discomfort, early satiety, postprandial fullness, bloating, or nausea. It may result from functional dyspepsia, gastroesophageal reflux, peptic ulcer disease, medications, or other gastrointestinal disorders.

\medterm{Infectious Esophagitis}-- Inflammation of the esophagus caused by an infectious organism, most often Candida, herpes simplex virus, or cytomegalovirus in immunocompromised patients. It commonly causes painful swallowing, difficulty swallowing, and retrosternal pain.

\medterm{Pancreas Divisum}-- A congenital anomaly in which the dorsal and ventral pancreatic ducts fail to fuse. Most people are asymptomatic, but impaired drainage through the minor papilla can predispose to recurrent pancreatitis.

\medterm{Pancreatic Abscess}-- A localized collection of pus within or adjacent to the pancreas, usually developing as a complication of infected pancreatic necrosis or severe pancreatitis. It may cause persistent fever, abdominal pain, sepsis, and organ failure.

\medterm{Parapneumonic Pleural Effusion}-- Pleural fluid accumulation associated with bacterial pneumonia. It may be uncomplicated or progress to a complicated effusion or empyema, requiring imaging, pleural-fluid analysis, antibiotics, and sometimes drainage.

\medterm{Pilonidal Sinus Disease}-- A chronic sinus or cyst in the natal cleft, usually containing hair and debris. It may cause recurrent pain, swelling, drainage, or abscess formation.

\medterm{Rumination Disorder}-- Recurrent effortless regurgitation of recently ingested food, followed by rechewing, reswallowing, or spitting, without the retching typical of vomiting. It can cause weight loss, malnutrition, and dental complications.

\medterm{Acalculous Cholecystitis}-- Acute inflammation of the gallbladder without gallstones, usually occurring in critically ill, postoperative, traumatised, or fasting patients. Presents with fever, right-upper-quadrant pain, leukocytosis, or unexplained sepsis; diagnosis is supported by ultrasound or cholescintigraphy. It can progress rapidly to gangrene or perforation, so treatment requires antibiotics and urgent surgical or percutaneous drainage.

\medterm{Aerophagy}-- Excessive swallowing of air, also called aerophagia, which can lead to frequent belching, abdominal bloating, distension, and discomfort. It may result from rapid eating or drinking, chewing gum, drinking carbonated beverages, talking while eating, mouth breathing, anxiety, or poorly coordinated swallowing; it can also occur with functional gastrointestinal disorders. Management involves identifying and reducing triggers, eating slowly, avoiding carbonated drinks and gum, and treating associated reflux, anxiety, or swallowing dysfunction.

\medterm{Diffuse Esophageal Spasm}-- A motility disorder characterised by premature, uncoordinated esophageal contractions, causing intermittent chest pain and dysphagia for liquids and solids. Endoscopy excludes structural disease; high-resolution manometry confirms premature contractions. Management may include dietary measures, smooth-muscle relaxants, and specialist treatment such as botulinum toxin or myotomy for refractory disease.

\medterm{Achlorhydria}-- Absence of hydrochloric acid in gastric secretions, defined as a gastric pH persistently $\ge$ 7. Causes: chronic autoimmune atrophic gastritis (pernicious anemia; anti-parietal cell and anti-intrinsic factor antibodies), chronic Helicobacter pylori gastritis, prolonged proton-pump inhibitor therapy, post-vagotomy, gastric carcinoma, and immunodeficiency states. Clinical consequences: increased susceptibility to enteric infections (Salmonella, Campylobacter, Vibrio), bacterial overgrowth, malabsorption of iron and vitamin B12, and hypergastrinemia (enterochromaffin-like cell hyperplasia; possible type II gastric carcinoid). Diagnosis: gastric pH monitoring, secretin-stimulation test where appropriate. Treatment: address underlying cause; parenteral vitamin B12 supplementation; surveillance endoscopy.

\medterm{Acute Pancreatitis}-- Acute inflammatory process of the pancreas most commonly caused by gallstones and alcohol accounting for approximately eighty percent of cases. The pathogenesis involves premature trypsin activation within pancreatic acinar cells triggering autodigestion. Clinical presentation includes severe epigastric pain radiating to the back, nausea, and vomiting. Diagnosis requires two of three criteria: characteristic pain, serum lipase at least three times upper limit of normal, and characteristic imaging findings of pancreatic inflammation on CT.

\medterm{Autoimmune Gastritis}-- Chronic atrophic gastritis caused by antibodies against parietal cells leading to achlorhydria and vitamin B12 deficiency. Associated with pernicious anemia and increased risk of gastric carcinoid tumors and type 1 gastric neuroendocrine tumors. Anti-parietal cell antibodies are present in approximately ninety percent of patients. Diagnosis is by serology showing elevated gastrin levels and decreased pepsinogen I to II ratio, confirmed by endoscopic biopsy showing oxyntic gland atrophy. Management involves iron supplementation and endoscopic surveillance for gastric neuroendocrine tumors.

\medterm{Autoimmune Enteropathy}-- Rare immune-mediated disease causing persistent diarrhoea and small-intestinal mucosal injury, usually with villous atrophy and crypt hyperplasia. It may occur in children or adults and can be associated with other autoimmune disorders or immune dysregulation. Symptoms include chronic diarrhoea, weight loss, abdominal distension, and nutritional deficiencies. Diagnosis requires compatible small-bowel biopsy findings, exclusion of infection and coeliac disease, and supportive serological or immunological findings where available.

\medterm{Barium Swallow}-- Fluoroscopic examination of the esophagus after ingestion of barium contrast. Used to evaluate esophageal anatomy, motility, strictures, and structural abnormalities. The study demonstrates mucosal detail, filling defects, strictures, and motility disorders. A barium follow-through extends the examination to the small bowel. Double-contrast technique using barium and air provides superior mucosal detail. The study has been largely replaced by upper endoscopy for direct visualization but remains useful in selected patients who cannot tolerate endoscopy.

\medterm{Bacterial Overgrowth}-- Excessive proliferation of bacteria in the small intestine, commonly called small-intestinal bacterial overgrowth. It may result from impaired intestinal motility, anatomical stasis, blind loops, strictures, diverticula, or loss of the normal ileocaecal barrier. Consequences include bloating, abdominal discomfort, diarrhoea, steatorrhoea, weight loss, and deficiencies of vitamin B12 and other nutrients. Diagnosis may involve breath testing or small-bowel aspirate culture.

\medterm{Barrett Esophagus}-- Metaplastic transformation of the normal stratified squamous epithelium of the distal esophagus to specialized intestinal-type columnar epithelium containing goblet cells, occurring as a consequence of chronic gastroesophageal reflux disease. This condition represents an important premalignant lesion because it confers a significantly increased risk of esophageal adenocarcinoma, estimated at approximately 0.5 percent per year in patients without dysplasia. The pathogenesis involves chronic acid exposure leading to inflammation, ulceration, and metaplastic change, with progression driven by bile reflux and inflammatory mediators.

\medterm{Barrett Esophagus Surveillance}-- Systematic endoscopic monitoring of patients with Barrett esophagus to detect dysplasia and early esophageal adenocarcinoma. Patients without dysplasia undergo surveillance every three to five years with four-quadrant biopsies every two centimeters throughout the Barrett segment. Patients with low-grade dysplasia should be offered endoscopic eradication therapy with radiofrequency ablation. Patients with high-grade dysplasia or intramucosal carcinoma should undergo endoscopic mucosal resection or oesophagectomy, depending on depth of invasion and patient fitness.

\medterm{Bile Reflux}-- Reflux of duodenal contents (bile acids, lysolecithin, pancreatic enzymes) into the stomach and esophagus, caused by incompetence of the pyloric sphincter, prior gastric surgery (Billroth II, Roux-en-Y, vagotomy with pyloroplasty), cholecystectomy, or functional gastroparesis. Causes bile reflux gastritis (chemical gastropathy): epigastric pain, nausea, bilious vomiting, weight loss; in the esophagus contributes to Barrett esophagus and refractory GERD. Diagnosis difficult; supports include endoscopic appearance of bile pooling, histologic chemical gastropathy (foveolar hyperplasia with vasodilation and edema), Bilitec (ambulatory bilirubin monitoring). Treatment: ursodeoxycholic acid, sucralfate, prokinetic agents, anti-reflux surgery (Roux-en-Y diversion) for severe cases.

\medterm{Bleeding Esophageal Varices}-- Life-threatening hemorrhage from dilated submucosal veins in the lower esophagus, occurring as a complication of portal hypertension, most commonly from cirrhosis. Varices form when the portal venous pressure gradient exceeds 10 mmHg (the threshold for variceal development). Risk of first variceal hemorrhage is approximately 15-25 percent per year, and each episode carries a mortality of 15-30 percent. Clinical presentation includes hematemesis (bright red blood or coffee-ground emesis), melena, and signs of hypovolaemic shock.

\medterm{Blind Loop Syndrome}-- Malabsorption syndrome due to stasis and bacterial overgrowth (SIBO) in a non-propulsive ("blind") segment of small bowel, classically caused by surgical bypass loops (Billroth II, end-to-side anastomosis), strictures (Crohn disease, radiation, adhesions), diverticulosis, or impaired motility (diabetes, scleroderma). Bacteria deconjugate bile salts and consume vitamin B12 producing steatorrhea, megaloblastic anemia, weight loss, diarrhoea, and malnutrition. Diagnosis: ${}^{14}$C-D-xylose or ${}^{14}$C-cholylglycine breath test, hydrogen breath test, response to empirical antibiotic course, jejunal aspirate culture $>10^5$ CFU/mL. Treatment: surgical correction of the blind loop, rotating broad-spectrum antibiotics (rifaximin, metronidazole, tetracycline, ciprofloxacin cycles), nutritional supplementation (vitamin B12, fat-soluble vitamins).

\medterm{Boerhaave Syndrome}-- Transmural spontaneous rupture of the esophagus typically occurring after forceful vomiting in the setting of a full stomach. Most commonly involves the left posterolateral wall of the distal esophagus. Presents with severe chest pain, subcutaneous emphysema, mediastinal emphysema, and rapidly progressive sepsis. This is a surgical emergency with mortality exceeding fifty percent when diagnosis and treatment are delayed. Diagnosis is suggested by chest radiography showing pneumomediastinum or pleural effusion on chest radiograph, and confirmed by CT showing esophageal wall thickening with extraluminal air.

\medterm{Bouveret Syndrome}-- Rare variant of gallstone ileus where a large gallstone erodes through a cholecystoduodenal fistula and impacts in the duodenum or pylorus causing gastric outlet obstruction. Represents approximately two to three percent of gallstone ileus cases. Presents with nausea, vomiting, epigastric pain, and anorexia. Diagnosis is by CT or upper gastrointestinal series showing gastric dilation with a filling defect in the duodenum and pneumobilia. Treatment is endoscopic or surgical removal of the impacted stone.

\medterm{Cameron Lesions}-- Linear gastric erosions or ulcers found on the mucosal folds at the diaphragmatic hiatus in patients with large hiatal hernias. They result from mechanical trauma as the gastric mucosa prolapses through the diaphragmatic crura during respiration and movement. May be a source of occult or overt gastrointestinal bleeding presenting with iron deficiency anemia or melena. Diagnosis requires careful endoscopic examination with particular attention to the hiatal hernia sac. Treatment involves proton pump inhibitors and treatment of the underlying mechanical cause.

\medterm{Caroli Disease}-- Congenital, autosomal recessive fibropolycystic liver disease (PKHD1 gene, encodes fibrocystin/polyductin; same gene as infantile polycystic kidney disease) characterized by non-obstructive, segmental, saccular dilation of the larger intrahepatic bile ducts. Two forms: the pure form (Caroli disease, duct dilation alone) and the more common combined form (Caroli syndrome, duct dilation + congenital hepatic fibrosis, presenting with portal hypertension, esophageal varices, hypersplenism). Often associated with medullary cystic kidney disease and autosomal recessive polycystic kidney disease. Presentation: recurrent episodes of cholangitis (RUQ pain, fever, jaundice), intrahepatic pigment gallstones, liver abscesses, biliary cirrhosis, increased risk of cholangiocarcinoma. Diagnosis: MRCP (diagnostic gold standard, avoids instrumentation of duct), cholangiography. Treatment: supportive — ursodeoxycholic acid for bile flow, antibiotics for cholangitis, endoscopic/surgical drainage of recurrent strictures and stones; liver transplantation for end-stage disease or refractory cholangitis.

\medterm{Cholestasis}-- Impairment of bile formation or flow leading to retention of bile acids, bilirubin, and cholesterol in the blood and tissues. Classified by site: hepatocellular (intrahepatic, eg, viral hepatitis, alcoholic hepatitis, drugs, primary biliary cholangitis, PSC, DILI, sepsis, pregnancy, hereditary) or obstructive (extrahepatic, eg, stones, strictures, tumors, biliary atresia, pancreatobiliary malignancy). Clinical features: prpruritus (worst at night, palms and soles), jaundice, dark urine, pale stools, fat malabsorption with steatorrhea and fat-soluble vitamin deficiency (A, D, E, K), xanthomas. Laboratory: elevated alkaline phosphatase, GGT, conjugated bilirubin, with mild AST/ALT rise (cholestatic pattern); high serum bile acids. Imaging (ultrasound, MRCP) localizes obstruction. Treatment: address underlying cause; bile acid sequestrants (cholestyramine), rifampicin, opioid antagonists (naltrexone) and sertraline for prpruritus; vitamin K and fat-soluble vitamin supplementation; biliary drainage (ERCP, surgery) for obstruction; liver transplant for end-stage disease.

\medterm{Cholelithiasis}-- Presence of gallstones in the gallbladder. Types: cholesterol stones (80\%, radiolucent, due to supersaturated bile), pigment stones (black: hemolysis, cirrhosis; brown: infection, bile stasis). Risk factors: female, fertile, fat, forty (5 Fs), obesity, rapid weight loss, hemolytic disorders, cirrhosis. Most are asymptomatic (incidental finding). Symptomatic cholelithiasis: episodic postprandial RUQ/epigastric pain (biliary colic). Complications: acute cholecystitis, choledocholithiasis, cholangitis, gallstone pancreatitis. Diagnosis: right upper quadrant ultrasound (gold standard). Treatment: laparoscopic cholecystectomy for symptomatic stones; observation for asymptomatic.

\medterm{Cushing Ulcer}-- Stress ulceration associated with central nervous system disease, trauma, or surgery. Caused by vagally mediated gastric acid hypersecretion, distinguishing it from Curling ulcer which results from mucosal ischemia. The ulcers are typically solitary, located in the stomach or duodenum, and present with gastrointestinal bleeding. Management involves acid suppression with proton pump inhibitors and treatment of the underlying neurological condition.

\medterm{Dieulafoy Lesion}-- Large caliber persistent submucosal artery in the stomach wall that can cause massive gastrointestinal bleeding. The artery is abnormally large and superficial, lacking the normal tapering into smaller arterioles. Presents with hematemesis or melena without a visible ulcer. The lesion is typically found within six centimeters of the gastroesophageal junction. Diagnosed and treated by endoscopic therapy with clips, thermal coagulation, or band ligation. Recurrent bleeding may occur and repeat endoscopy with hemostatic therapy.

\medterm{Digestion}-- The process by which the body breaks down ingested food into smaller molecules that can be absorbed and utilized. Digestion begins in the mouth (mechanical--mastication; enzymatic--salivary amylase for starch, lingual lipase). In the stomach: gastric acid (HCl) denatures proteins, pepsin digests proteins, and mechanical churning produces chyme. In the small intestine: pancreatic enzymes (trypsin, chymotrypsin, amylase, lipase, nucleases) break down macromolecules; bile salts emulsify fats. Brush border enzymes (lactase, sucrase, maltase, peptidases) complete digestion. Absorption occurs primarily in the small intestine. The colon absorbs water and electrolytes and ferments undigested fiber.

\medterm{Digestive Health}-- The efficient functioning of the gastrointestinal tract in breaking down food, absorbing nutrients, and eliminating waste. Digestive health depends on factors such as diet, hydration, gut microbiome balance, and the integrity of the digestive mucosa. Common digestive disorders include irritable bowel syndrome, gastroesophageal reflux disease, and inflammatory bowel disease.

\medterm{Dumping Syndrome}-- Rapid gastric emptying of hypertonic contents into the small bowel after gastrectomy or gastric bypass surgery. Early dumping occurs within thirty minutes with vasomotor symptoms including flushing, palpitations, diarrhea, and dizziness from fluid shifts into the intestinal lumen. Late dumping occurs one to three hours later with reactive hypoglycemia from exaggerated insulin secretion. Managed with dietary modifications including small frequent meals, avoidance of simple sugars, and lying down after meals. Pharmacological therapy includes octreotide and acarbose to slow gastric emptying and blunt the insulin response.

\medterm{Duodenal Ulcer}-- Peptic ulcer in the duodenal mucosa strongly associated with H. pylori infection. Presents with epigastric pain relieved by food typically occurring several hours after meals and at night. Complications include hemorrhage, perforation, and gastric outlet obstruction. Diagnosis is by upper endoscopy. Treatment involves proton pump inhibitor therapy for four to eight weeks and H. pylori eradication. Duodenal ulcers have a lower risk of malignancy compared to gastric ulcers but require follow-up endoscopy to document healing.

\medterm{Duodenitis}-- Inflammation of the duodenal mucosa, caused by Helicobacter pylori infection, acid-related disease, nonsteroidal anti-inflammatory drugs, infection, coeliac disease, or inflammatory bowel disease. Symptoms may include epigastric pain, nausea, vomiting, and gastrointestinal bleeding. Diagnosis is usually by upper endoscopy with biopsy, and treatment is directed at the underlying cause.

\medterm{Dyspepsia}-- Chronic or recurrent pain or discomfort centered in the upper abdomen. Organic causes: peptic ulcer disease, GERD, gastritis, gastric cancer. Functional dyspepsia is diagnosed after exclusion of organic pathology with normal upper endoscopy. Rome IV criteria define subtypes: epigastric pain syndrome and postprandial distress syndrome. Management: H. pylori test and treat, PPI trial, prokinetics, and low-dose antidepressants.

\medterm{Dysphagia}-- Difficulty swallowing, classified as oropharyngeal (difficulty initiating a swallow, choking, coughing, nasal regurgitation) or esophageal (sensation of food sticking retrosternally). Oropharyngeal causes: neurological (stroke, Parkinson's disease, motor neurone disease, myasthenia gravis), structural (pharyngeal pouch, tumors), and muscular (dermatomyositis). esophageal causes: mechanical obstruction (stricture, Schatzki ring, esophageal cancer, eosinophilic esophagitis) and motility disorders (achalasia, diffuse esophageal spasm, scleroderma). Investigations: videofluoroscopic swallow study for oropharyngeal; upper endoscopy, barium swallow, and esophageal manometry for esophageal. Complications: aspiration pneumonia, malnutrition, and dehydration.

\medterm{EGD}-- Esophagogastroduodenoscopy: upper endoscopy visualizing the esophagus, stomach, and duodenum using a flexible endoscope. Indications: dyspepsia with alarm features, dysphagia, GI bleeding (diagnosis and therapy—injection, clips, banding, cautery), biopsy (celiac, H. pylori, Barrett surveillance), dilation of strictures, and placement of feeding tubes. Complications: perforation, bleeding, sedation-related events. Preparation: fasting; sedation usually light. Alternatives for screening: capsule endoscopy. Essential for diagnosis and treatment of upper GI disorders.

\medterm{Emesis (Vomiting)}-- The forceful expulsion of gastric contents through the mouth, a complex reflex coordinated by the vomiting center in the medulla oblongata. The reflex involves: prodromal phase (nausea, salivation, autonomic activation--tachycardia, sweating), retching (rhythmic contractions of diaphragm and abdominal muscles against a closed glottis), and expulsion (forceful contraction of abdominal muscles with relaxation of the lower esophageal sphincter and diaphragm). Causes: gastrointestinal (gastroenteritis, gastroparesis, obstruction), vestibular (motion sickness, labyrinthitis), CNS (increased intracranial pressure, migraine), metabolic (pregnancy, uremia, DKA, hypercalcemia), medications (chemotherapy, opioids, antibiotics), toxins (food poisoning, alcohol), and psychogenic. Complications: aspiration, dehydration, electrolyte disturbances, Mallory-Weiss tear. Antiemetics target different receptors: 5-HT3, D2, H1, M1, NK1.

\medterm{Endoscopic Mucosal Resection}-- Advanced endoscopic technique for removal of superficial gastrointestinal neoplasms by lifting the lesion through submucosal injection and resecting with an electrosurgical snare. The technique involves injection of saline with vasoconstrictive agent into the submucosa, capturing the elevated lesion with a snare, and resecting with electrocautery. Indications include sessile colorectal polyps, early gastric cancer, Barrett esophagus with high-grade dysplasia, and large duodenal lesions. Complications include bleeding, perforation, and stricture formation, but these are uncommon in experienced hands.

\medterm{Endoscopic Submucosal Dissection}-- Advanced endoscopic technique for en bloc resection of large gastrointestinal neoplasms using specialized electrosurgical knives to dissect through the submucosa. The technique allows complete histological assessment and is particularly valuable for early gastric and colorectal cancers. Unlike piecemeal EMR, ESD achieves higher en bloc resection rates, lower recurrence rates, and accurate staging. Complications include bleeding in five to ten percent and perforation in three to five percent, though the risk of perforation is higher than with EMR and requires careful patient selection and operator expertise.

\medterm{Endoscopic Ultrasound}-- Diagnostic and therapeutic endoscopic procedure combining endoscopy with high-frequency ultrasound to provide detailed imaging of the gastrointestinal wall layers and adjacent organs. EUS is particularly valuable for evaluating pancreatic masses, staging gastrointestinal malignancies, evaluating submucosal lesions, and assessing the biliary tree. EUS-guided fine-needle aspiration allows tissue sampling of lesions not accessible by conventional biopsy. The staging accuracy for esophageal and rectal cancers, and for evaluating pancreatic lesions and mediastinal lymphadenopathy.

\medterm{Endoscopy}-- Visualization of the gastrointestinal tract using flexible fiber-optic or video endoscopes. Upper endoscopy examines the esophagus, stomach, and duodenum. Colonoscopy examines the colon and terminal ileum. Both allow direct visualization, biopsy, and therapeutic interventions including polypectomy, hemostasis, and dilation. Endoscopic procedures require patient preparation, sedation, and monitoring. Complications are uncommon but include perforation, bleeding, and adverse reactions to sedation. Most procedures are safe when performed by trained endoscopists with appropriate monitoring.

\medterm{Eosinophilic Esophagitis}-- Chronic immune-mediated condition characterized by dense eosinophilic infiltration of the esophageal mucosa defined as fifteen or more eosinophils per high-power field, in the absence of other causes of esophageal eosinophilia. The condition is driven by a T helper type two inflammatory response and is strongly associated with atopic conditions. The clinical presentation varies by age: infants present with feeding difficulties; school-age children with dysphagia; and adults with progressive dysphagia, and food bolus impaction.

\medterm{Eosinophilic Gastroenteritis}-- Rare chronic inflammatory condition characterized by eosinophilic infiltration of the gastrointestinal tract beyond the esophagus involving any segment from stomach to rectum and extending beyond the mucosa to involve the muscularis propria or serosa. The condition is classified based on depth of infiltration: mucosal disease presents with pain and diarrhea; muscular disease causes obstruction; serosal disease causes eosinophilic ascites. Diagnosis is established by tissue biopsy demonstrating eosinophilic infiltration in the affected tissue.

\medterm{Erosive Gastritis}-- Inflammation of the stomach lining with superficial erosions but no full-thickness ulceration. Commonly caused by NSAIDs, alcohol, or critical illness. May present with epigastric pain or gastrointestinal bleeding. Diagnosis is by upper endoscopy showing erosions and erythema. Management involves proton pump inhibitors and removal of offending agents. Stress ulcer prophylaxis with proton pump inhibitors or H2 blockers is indicated in critically ill patients at risk for stress-related mucosal disease.

\medterm{Esophageal Achalasia}-- Primary esophageal motility disorder characterized by failure of the lower esophageal sphincter to relax during swallowing and absent peristalsis in the distal esophagus, resulting from selective loss of inhibitory neurons in the myenteric plexus. The condition typically presents with progressive dysphagia to both solids and liquids, regurgitation, chest pain, and weight loss. Complications include aspiration pneumonia and increased risk of esophageal squamous cell carcinoma. Diagnosis is confirmed by high-resolution manometry showing incomplete lower esophageal sphincter relaxation and absent peristalsis.

\medterm{Esophageal Candidiasis}-- Fungal infection of the esophagus caused by Candida species common in immunocompromised patients including those with HIV, organ transplant recipients, and patients on inhaled corticosteroids. Presents with odynophagia and dysphagia. Diagnosed by endoscopy showing white adherent plaques and biopsy confirming fungal hyphae. Treated with oral fluconazole as first-line therapy. Esophageal candidiasis is an AIDS-defining illness when the CD4 count is below two hundred cells per microliter.

\medterm{Esophageal Dilation}-- Endoscopic passage of bougies or balloons to widen esophageal strictures. Used for symptomatic strictures from peptic disease, radiation, caustic injury, or eosinophilic esophagitis. Savary bougies are passed over a guidewire with progressive dilation. Through-the-scope balloons allow controlled radial dilation under direct visualization. Complications include perforation, which occurs in approximately one percent, and bleeding. Serial dilations are often required for refractory strictures. Dilation should be performed gradually over multiple sessions to reduce perforation risk.

\medterm{Esophageal Diverticulum}-- Outpouching of the esophageal wall classified by location. Zenker diverticulum occurs in the hypopharynx and is the most common esophageal diverticulum. Mid-esophageal diverticula are traction diverticula caused by extrinsic inflammatory disease including tuberculosis and histoplasmosis. Epiphrenic diverticula occur in the distal esophagus above the lower esophageal sphincter and are associated with esophageal motility disorders including achalasia and diffuse esophageal spasm. Most diverticula are asymptomatic and discovered incidentally; symptomatic cases may require surgical diverticulectomy and myotomy.

\medterm{Esophageal Intramural Pseudodiverticulosis}-- Condition characterized by multiple small outpouchings of the esophageal submucosal glands seen on barium swallow as flask-shaped contrast collections oriented longitudinally along the esophagus. The pseudodiverticula are not true diverticula but rather dilated excretory ducts of submucosal glands. Associated with gastroesophageal reflux disease, diabetes mellitus, and chronic alcoholism. Most patients present with progressive dysphagia due to associated esophageal stricture formation. Diagnosis is confirmed by barium swallow showing characteristic flask-shaped outpouchings, and management focuses on treating the underlying stricture with dilation.

\medterm{Esophageal Leukoplakia}-- White plaques on the esophageal mucosa representing squamous cell hyperplasia with hyperkeratosis. Similar to oral leukoplakia it is considered a potentially premalignant condition. Associated with chronic irritation from smoking, alcohol, and gastroesophageal reflux. Diagnosis is by endoscopy with biopsy to exclude dysplasia or carcinoma. Management includes surveillance with endoscopy and biopsy, smoking and alcohol cessation, and treatment of reflux disease. The risk of malignant transformation is low, but regular surveillance with endoscopic biopsy is recommended.

\medterm{Esophageal Manometry}-- Diagnostic test measuring esophageal motor function by recording intraluminal pressures during swallowing. A catheter with pressure sensors is passed through the nose into the stomach and withdrawn to record pressures at multiple levels. High-resolution manometry using closely spaced sensors provides a spatiotemporal color plot. The Chicago classification categorizes esophageal motility disorders into major disorders including achalasia and distal esophageal spasm, and minor disorders including ineffective esophageal motility (IEM).

\medterm{Esophageal Perforation}-- Full-thickness rupture of the esophageal wall, most commonly Boerhaave syndrome from forceful vomiting occurring on the left posterolateral distal esophagus. Iatrogenic perforation during endoscopy is another common cause. Presents with chest pain, subcutaneous emphysema, and mediastinitis. A surgical emergency requiring emergent repair and drainage. Delayed diagnosis significantly increases mortality. Management includes nil per os, intravenous antibiotics, and surgical repair with tissue flap reinforcement for large defects.

\medterm{Esophageal Spasm}-- Esophageal motility disorder characterized by uncoordinated or excessive contractions impairing bolus transport. Diffuse esophageal spasm shows simultaneous contractions on manometry. Jackhammer esophagus features extremely high-amplitude peristaltic contractions. Patients present with episodic dysphagia and chest pain that may mimic cardiac pain. Diagnosis is by high-resolution manometry demonstrating characteristic pressure patterns. Barium swallow may show the corkscrew esophagus pattern. Treatment includes calcium channel blockers, nitrates, phosphodiesterase inhibitors, and endoscopic botulinum toxin injection for symptomatic relief.

\medterm{Esophageal Stricture}-- Narrowing of the esophageal lumen from chronic gastroesophageal reflux disease, caustic ingestion, or eosinophilic esophagitis. Presents with progressive dysphagia initially to solids then liquids. Diagnosed by upper endoscopy and barium swallow. Peptic strictures are the most common type and are treated with endoscopic balloon dilation combined with aggressive acid suppression. Recurrent strictures may require repeated dilations. Malignant strictures require stent placement or surgical resection with reconstruction.

\medterm{Esophageal Varices}-- Dilated, tortuous submucosal veins in the lower esophagus, most commonly caused by portal hypertension (cirrhosis—hepatic fibrosis with increased portal venous pressure >10 mmHg). They develop as portosystemic collaterals connecting the portal venous system (left gastric vein) to the systemic venous system (azygos vein). Prevalence: ~50\% of patients with cirrhosis, 30\% of whom will bleed within 2 years. Ruptured esophageal varices cause acute upper GI hemorrhage (hematemesis, melena, hemodynamic instability), with mortality of 15--20\% per episode. Diagnosis: upper endoscopy—emergency for suspected variceal bleeding (red wale signs, cherry-red spots, fibrin plug). Management of acute bleeding: resuscitation, vasoactive drugs (terlipressin or octreotide), prophylactic antibiotics (ceftriaxone 1 g IV, reduces bacterial peritonitis and mortality), endoscopic band ligation (first-line, 90\% hemostasis), and covered transjugular intrahepatic portosystemic shunt (TIPS) for refractory bleeding. Secondary prophylaxis: non-selective beta-blockers (propranolol, carvedilol) plus repeat band ligation every 2--4 weeks until varices obliterated. Primary prophylaxis for medium/large varices: beta-blockers or band ligation.

\medterm{Esophageal Web}-- Thin mucosal fold protruding into the esophageal lumen most commonly in the postcricoid region or upper esophagus. Webs may be asymptomatic or cause dysphraga particularly to solids. Multiple webs can occur in association with iron deficiency anemia as part of Plummer-Vinson syndrome or Paterson-Kelly syndrome. Diagnosis is by barium swallow showing thin filling defects and upper endoscopy which may inadvertently displace the web. Symptomatic webs are treated with endoscopic dilation. The significance of an esophageal web lies primarily in its association with Plummer-Vinson syndrome and the increased risk of postcricoid squamous cell carcinoma.

\medterm{Esophagitis}-- Inflammation of the esophageal mucosa. Causes: gastro-esophageal reflux disease (GERD—most common, ~60--70\% of cases), infectious (Candida albicans—most common in immunocompromised, HIV, diabetes, post-antibiotics; herpes simplex virus—painful ulcers; CMV—large ulcerations in immunocompromised), eosinophilic esophagitis (EoE—allergic/atopic, characteristic rings, furrows, and white plaques on endoscopy), pill-induced (doxycycline, bisphosphonates, NSAIDs, potassium chloride, quinidine—typically in the mid-esophagus at areas of extrinsic compression), and radiation esophagitis (during thoracic radiotherapy). Symptoms: heartburn, odynophagia (painful swallowing), dysphagia (especially for solids in EoE and strictures), and retrosternal chest discomfort. Diagnosis: upper endoscopy with biopsy is the gold standard; biopsies from the distal and mid-esophagus help distinguish reflux esophagitis (basal hyperplasia, eosinophils >15/hpf in EoE, neutrophilic infiltration in severe reflux). GERD esophagitis is graded by the Los Angeles classification (A--D). Management: lifestyle measures plus proton pump inhibitors (omeprazole 20 mg, lansoprazole 30 mg daily) for 4--8 weeks; EoE requires topical swallowed corticosteroids (fluticasone, budesonide) and dietary elimination; infectious esophagitis is treated with antifungals or antivirals as appropriate.

\medterm{Esophagus (esophagus)}-- A muscular tube approximately 25 cm long connecting the pharynx to the stomach, extending from the cricopharyngeus muscle (at the level of C6) to the gastroesophageal junction (T11). The esophagus has three anatomic constrictions: cervical (upper esophageal sphincter--cricopharyngeus at C6), thoracic (crossing of the aortic arch and left main bronchus), and diaphragmatic (hiatus). The wall consists of mucosa (non-keratinized stratified squamous epithelium), submucosa, muscularis propria (upper third: skeletal muscle; middle third: mixed; lower third: smooth muscle), and adventitia (no serosa). Functions: transport of swallowed food and liquid from pharynx to stomach via coordinated peristaltic contractions. The lower esophageal sphincter prevents gastric reflux. Common disorders: GERD, esophagitis, Barrett esophagus, esophageal cancer, achalasia, and strictures.

\medterm{Feeding Jejunostomy}-- Surgical or endoscopic placement of a feeding tube into the jejunum for enteral nutrition when upper gastrointestinal access is not feasible. Used in patients with head and neck cancers, esophageal obstruction, or severe dysphagia. Provides long-term nutritional support and avoids the complications of parenteral nutrition. The tube can be placed percutaneously using endoscopic or radiological guidance. Complications include tube displacement, occlusion, leakage, and local infection.

\medterm{Functional Dyspepsia}-- Chronic epigastric pain or discomfort without identifiable structural cause. The Rome IV criteria include postprandial distress syndrome with early satiety and postprandial fullness, and epigastric pain syndrome with epigastric pain or burning unrelated to meals. Pathophysiology involves visceral hypersensitivity, altered gastric accommodation, and duodenal eosinophilia. Management includes dietary modifications, proton pump inhibitors, histamine H2 receptor antagonists, prokinetics, and neuromodulators such as tricyclic antidepressants and serotonin-norepinephrine reuptake inhibitors.

\medterm{Gastralgia}-- Pain in the stomach or epigastric region. The term is non-specific and encompasses various causes of epigastric discomfort: functional dyspepsia (most common), peptic ulcer disease (gastric or duodenal), gastritis, gastro-esophageal reflux disease (GERD), gastric cancer, pancreatitis, cholecystitis, and referred pain (myocardial ischemia, pleurisy). Clinical assessment: quality (burning, gnawing, cramping), timing (postprandial, nocturnal, fasting), aggravating/relieving factors, associated symptoms (nausea, vomiting, weight loss, hematemesis, melena). Management: identify and treat underlying cause; H. pylori testing and eradication if present; trial of PPI if GERD or peptic ulcer suspected.

\medterm{Gastrectomy}-- Surgical removal of all (total gastrectomy) or part (partial/subtotal gastrectomy) of the stomach. Indications: gastric cancer (most common), severe peptic ulcer disease (rarely needed now due to PPIs and H. pylori treatment), gastric lymphoma, gastrointestinal stromal tumor (GIST), perforation, bleeding. Types: Billroth I (partial gastrectomy with gastroduodenostomy), Billroth II (partial gastrectomy with gastrojejunostomy), Roux-en-Y (total gastrectomy with oesophagojejunostomy). Complications: anastomotic leak (highest risk), dumping syndrome, bile reflux gastritis (post Billroth II), malabsorption (iron, B12, calcium, vitamin D, fat-soluble vitamins), bacterial overgrowth, weight loss, afferent/efferent loop syndrome. After total gastrectomy, intramuscular vitamin B12 supplementation is required lifelong due to loss of intrinsic factor.

\medterm{Gastric Bezoar}-- Accumulation of indigestible material in the stomach. Phytobezoars consist of vegetable fiber, trichobezoars of hair, and pharmacobezoars of medications. Associated with gastroparesis, prior surgery, and impaired gastric motility. Diagnosed by endoscopy. Treated with enzymatic dissolution using cellulase, papain, or cola, endoscopic fragmentation, or surgical removal for refractory cases. Bezoars are an important cause of gastric outlet obstruction and ulceration.

\medterm{Gastric Carcinoid Tumor}-- Neuroendocrine tumor arising from enterochromaffin-like cells in the gastric fundus. Type I gastric carcinoids are associated with autoimmune atrophic gastritis and hypergastrinemia, are typically small and multifocal, and carry a good prognosis. Type II gastric carcinoids are associated with Zollinger-Ellison syndrome and multiple endocrine neoplasia type 1. Type III gastric carcinoids are sporadic, solitary, larger, and more aggressive. Diagnosed by endoscopy with biopsy. Treatment depends on tumor size, type, and stage: small type I tumors can be managed by endoscopic resection and surveillance, while types II and III often require surgical resection with lymphadenectomy.

\medterm{Gastric Emptying Study}-- Nuclear medicine diagnostic test measuring the rate of gastric emptying. The protocol involves ingestion of a radiolabeled meal followed by serial gamma camera images obtained at one, two, and four hours. Gastric emptying is considered delayed if more than ten percent of the meal remains at four hours. The test is primarily used to diagnose gastroparesis, a condition of delayed gastric emptying without mechanical obstruction most commonly occurring as a complication of diabetes mellitus. The study may also help differentiate gastroparesis from functional dyspepsia.

\medterm{Gastric Lymphoma}-- Malignant proliferation of lymphoid tissue in the stomach most commonly mucosa-associated lymphoid tissue lymphoma. Strongly associated with H. pylori infection. Early disease may regress completely with H. pylori eradication therapy alone. More advanced disease may require chemotherapy with rituximab and cyclophosphamide, doxorubicin, vincristine, and prednisone. Surgery is rarely needed. Diffuse large B-cell lymphoma of the stomach is the other common type and requires combination chemotherapy with R-CHOP (rituximab, cyclophosphamide, doxorubicin, vincristine, prednisolone).

\medterm{Gastric Outlet Obstruction}-- Blockage of gastric outflow at the pylorus or duodenum. Causes include peptic ulcer scarring, gastric cancer, and pyloric stenosis in infants. Presents with projectile vomiting, dehydration, and metabolic alkalosis. Diagnosis is by upper endoscopy and CT scan. Management requires decompression with nasogastric tube, fluid resuscitation, correction of electrolyte abnormalities, and definitive treatment of the underlying cause. Pyloric stenosis in infants is treated surgically with pyloromyotomy.

\medterm{Gastric Polyps}-- Protruding lesions of the gastric mucosa classified by histological type. Fundic gland polyps are the most common type associated with long-term proton pump inhibitor use and familial adenomatous polyposis. Hyperplastic polyps are associated with chronic gastritis and H. pylori infection and carry a small risk of dysplasia. Adenomatous polyps are true neoplasms with malignant potential requiring complete removal and surveillance. Inflammatory polyps occur in the setting of inflammatory bowel disease and may regress with treatment of the underlying colitis.

\medterm{Gastric Ulcer}-- Peptic ulcer in the gastric mucosa associated with H. pylori, NSAIDs, and smoking. Pain is typically worsened by food. There is a higher risk of malignancy than duodenal ulcers. Diagnosis is by upper endoscopy with biopsy to exclude malignancy. Treatment involves proton pump inhibitors and H. pylori eradication. Follow-up endoscopy at eight weeks is recommended to confirm healing and exclude malignancy in gastric ulcers, particularly in patients over forty-five years.

\medterm{Gastric Ulcer (see Peptic Ulcer)}-- See Peptic Ulcer.

\medterm{Gastric Volvulus}-- Rotation of the stomach along its axis causing obstruction. It may present with Borchardt triad of epigastric pain, retching, and inability to pass a nasogastric tube. It is a surgical emergency requiring decompression and gastropexy. Organoaxial volvulus is the most common type where the stomach rotates along its longitudinal axis. Mesenteroaxial volvulus involves rotation along an axis perpendicular to the long axis. Diagnosis is by CT scan or upper endoscopy.

\medterm{Gastrin}-- A peptide hormone produced by G cells in the gastric antrum and duodenum. Secreted in response to: (1) gastric distension, (2) amino acids and peptides (especially from protein-rich meals), (3) vagal stimulation (via gastrin-releasing peptide, GRP), (4) hypercalcaemia. Actions: stimulates gastric acid secretion by parietal cells (via histamine release from enterochromaffin-like cells and direct effect on CCK-B receptors), stimulates gastric motility, promotes growth of gastric mucosa (trophic effect). Inhibited by somatostatin (from D cells in response to low antral pH). Hypergastrinaemia: Zollinger-Ellison syndrome (gastrinoma, usually pancreatic), pernicious anemia (achlorhydria, loss of feedback inhibition), chronic PPI use. Fasting serum gastrin is measured in suspected gastrinoma.

\medterm{Gastritis}-- Inflammation of the gastric mucosa, classified by duration (acute or chronic), etiology, and histologic pattern. Acute gastritis is most commonly caused by Helicobacter pylori infection (acute H. pylori infection causes transient hypochlorhydria with epigastric pain, nausea, and vomiting), NSAIDs (inhibiting mucosal prostaglandin synthesis, especially in the gastric antrum, leading to erosions and bleeding), excessive alcohol consumption, severe physiologic stress (stress-related gastritis in critically ill patients). Chronic gastritis is most commonly caused by persistent H. pylori infection, which can progress to atrophic gastritis, intestinal metaplasia, and gastric adenocarcinoma.

\medterm{Gastroduodenal Crohn Disease}-- Involvement of the stomach and duodenum by Crohn disease occurring in approximately five percent of patients with Crohn. Gastric Crohn typically involves the antrum and pylorus and can cause gastric outlet obstruction from edema, fibrosis, or stricture. Duodenal Crohn presents with epigastric pain, nausea, vomiting, and weight loss. Diagnosis is by upper endoscopy with biopsy showing non-caseating granulomas and chronic inflammation. Treatment involves medical therapy with corticosteroids for induction of remission, followed by immunomodulators or biologic therapy for maintenance.

\medterm{Gastroenterology}-- The medical specialty focused on the diagnosis and management of diseases of the gastrointestinal tract (esophagus, stomach, small intestine, colon, rectum, pancreas, gallbladder, bile ducts, liver). Subspecialties: hepatology (liver), pancreatology, inflammatory bowel disease, gastrointestinal oncology, neurogastroenterology. Diagnostic tools: upper endoscopy (esophagogastroduodenoscopy, OGD), colonoscopy, capsule endoscopy, endoscopic retrograde cholangiopancreatography (ERCP), endoscopic ultrasound (EUS), breath tests (H. pylori, lactose intolerance, SIBO), manometry, pH metry. Therapeutic procedures: polypectomy, mucosal resection, hemostasis (injection, cautery, clips, banding), dilation, stent placement, percutaneous endoscopic gastrostomy (PEG).

\medterm{Gastroesophageal Reflux Disease (GERD)}-- A chronic condition in which gastric contents reflux backward into the esophagus, causing troublesome symptoms and/or complications. Epidemiology: affects 10--20\% of adults in Western countries. Pathophysiology: transient lower esophageal sphincter relaxations (TLESRs) are the primary mechanism; contributing factors include a hypotensive lower esophageal sphincter, hiatal hernia, impaired esophageal clearance, delayed gastric emptying, and obesity. Symptoms: classic symptoms include heartburn (retrosternal burning) and acid regurgitation; atypical symptoms include non-cardiac chest pain, chronic cough, laryngitis, asthma exacerbation, and globus sensation. Complications: erosive esophagitis, esophageal stricture, Barrett esophagus (metaplasia from squamous to columnar epithelium, increasing risk of esophageal adenocarcinoma). Diagnosis: clinical evaluation, PPI trial, upper endoscopy (for alarm symptoms such as dysphagia, weight loss, GI bleeding, anemia, or age > 55), and 24-hour ambulatory pH/impedance monitoring for refractory cases. Management: lifestyle modifications (weight loss, elevation of head of bed, avoiding trigger foods and late meals), proton pump inhibitors (PPIs like omeprazole 20 mg daily), H2 receptor antagonists, and laparoscopic Nissen fundoplication or magnetic sphincter augmentation for refractory cases.

\medterm{Gastrointestinal Bleeding}-- Hemorrhage from any part of the gastrointestinal tract classified as upper or lower based on location relative to the ligament of Treitz. Upper gastrointestinal bleeding originating proximal to this landmark typically presents with hematemesis or melena. Common causes include peptic ulcer disease, esophageal and gastric varices from portal hypertension, Mallory-Weiss tears, and erosive gastritis. Lower gastrointestinal bleeding originating distal to the ligament of Treitz usually presents with haematochezia or bright red blood per rectum, along with abdominal pain.

\medterm{Gastric Malignancy}-- See \hyperlink{Stomach Cancer}{stomach cancer}.

\medterm{Gastroparesis}-- Delayed gastric emptying without mechanical obstruction most commonly caused by diabetes mellitus, post-surgical changes, or idiopathic causes. Presents with nausea, vomiting of undigested food, early satiety, bloating, and abdominal pain. Diagnosed by gastric emptying scintigraphy showing more than ten percent retention at four hours. Complications include poor glycemic control, bezoar formation, and malnutrition. Managed with dietary modifications including small frequent meals and low-fat low-fibre diet, prokinetic agents such as metoclopramide and domperidone, and antiemetics. In severe cases, gastric electrical stimulation or jejunostomy tube placement may be considered.

\medterm{Gastrostomy}-- Surgical creation of an opening (stoma) through the abdominal wall into the stomach for enteral feeding or decompression. Methods: percutaneous endoscopic gastrostomy (PEG, most common--pull or introducer technique), radiologically inserted gastrostomy (RIG, under fluoroscopic guidance), surgical gastrostomy (open or laparoscopic; Stamm or Witzel technique). Indications: neurological dysphagia (stroke, motor neuron disease, dementia), head and neck cancer, esophageal obstruction, major trauma/burns, prolonged mechanical ventilation. Contraindications: gastric outlet obstruction, severe ascites, gastric varices, coagulopathy, previous abdominal surgery (relative). Complications: peristomal infection (most common, 5-30\%), tube displacement, buried bumper syndrome, gastrocolic fistula, peritonitis, pneumoperitoneum, bleeding.

\medterm{Gavage}-- Feeding through a tube passed into the stomach (nasogastric, orogastric, or gastrostomy tube). Used when oral feeding is impossible or inadequate. In neonates: gavage feeding is used for premature infants unable to coordinate suck-swallow-breathe. Types: bolus (intermittent, 100-400 mL over 10-30 minutes) or continuous (drip, 50-150 mL/hour via pump). Complications: tube misplacement (aspiration), tube blockage, nasal erosion, sinusitis, gastro-esophageal reflux, diarrhea (hyperosmolar feeds).

\medterm{Halitosis}-- An unpleasant odour of the breath that may be socially distressing. The most common cause (85-90\%) is poor oral hygiene-- accumulation of bacteria on the tongue, dental plaque, periodontitis, and gingivitis. The odour arises from volatile sulphur compounds (VSCs--hydrogen sulphide, methyl mercaptan, dimethyl sulphide) produced by Gram-negative anaerobic bacteria (especially Porphyromonas gingivalis, Treponema denticola, Tannerella forsythia) that break down sulphur-containing amino acids (cysteine, methionine). Other oral causes: dry mouth (xerostomia, reduced salivary cleansing), dental caries, oral infections, tongue coating, poorly fitting dentures. Non-oral causes: sinusitis (post-nasal drip), tonsillitis/tonsilloliths, gastro-esophageal reflux disease (GERD), Helicobacter pylori infection, liver failure (fetor hepaticus--a sweet, musty odour), renal failure (uraemic fetor--ammoniacal/urine-like), diabetic ketoacidosis (fruity/acetone), lung abscess, bronchiectasis (putrid odour), and trimethylaminuria (fish malodour syndrome). Diagnosis: clinical assessment (oral examination), organoleptic scoring (sensory assessment by clinician), gas chromatography, BANA test (benzoyl-DL-arginine-naphthylamide for bacterial enzymes). Management: good oral hygiene (tongue scraping, toothbrushing, flossing), treatment of underlying dental/medical conditions, mouthwash (chlorhexidine, zinc-containing, chlorine dioxide), avoidance of tobacco and alcohol.

\medterm{Helicobacter pylori}-- Gram-negative, microaerophilic, spiral-shaped bacterium specifically adapted to colonize the human gastric mucosa. It is estimated that approximately half the world's population is infected, making it one of the most common chronic bacterial infections globally. H. pylori colonizes the gastric antrum and body, surviving the extremely acidic environment by producing large quantities of the enzyme urease, which converts urea to ammonia and carbon dioxide, creating a protective alkaline microenvironment around the organism that neutralizes gastric acid and allows colonization.

\medterm{Hematemesis}-- Vomiting of blood, indicating upper gastrointestinal bleeding proximal to the ligament of Treitz. Blood may be bright red (acute, brisk bleeding) or resemble coffee grounds (altered blood that has been in contact with gastric acid). Common causes: peptic ulcer disease (most common, ~50\%), esophageal varices (~15\%), Mallory-Weiss tears (~10\%), erosive gastritis, esophagitis, and gastric cancer. Initial management: ABC approach, large-bore IV access, fluid resuscitation, blood transfusion (target Hb >80 g/L in hemodynamically stable patients), and proton pump inhibitor infusion. Prognostic scoring: Rockall score and Glasgow-Blatchford score predict need for intervention and mortality. Upper endoscopy within 24 hours is both diagnostic and therapeutic (hemostatic clips, adrenaline injection, thermal coagulation, or band ligation for varices).

\medterm{Hematochezia}-- The passage of bright red or maroon blood per rectum, usually indicating active lower gastrointestinal (LGI) bleeding distal to the ligament of Treitz (e.g., diverticulosis, arteriovenous malformations/angiodysplasia, colorectal polyps/cancer, inflammatory bowel disease, hemorrhoids). Massive upper GI bleeding (e.g., esophageal varices, peptic ulcer disease) with rapid intestinal transit ($<2$ hours) can also present with hematochezia and hemodynamic instability. Diagnostic evaluation: colonoscopy (first-line), CTA, or red blood cell scan.

\medterm{Hiatal Hernia}-- Protrusion of part of the stomach through the esophageal hiatus into the thoracic cavity. Type I (sliding, 95\%): gastroesophageal junction above the diaphragm. Types II-IV (paraesophageal): fundus or other portions herniate alongside the esophagus. Sliding hernias are associated with GERD. Paraesophageal hernias can cause obstruction, volvulus, and ischemia. Treatment: PPIs for GERD; surgical repair for symptomatic paraesophageal hernias.

\medterm{Hiccough (Hiccup)}-- An involuntary, repetitive, spasmodic contraction of the diaphragm followed by sudden closure of the glottis, producing the characteristic 'hic' sound. Hiccups are caused by irritation or stimulation of the phrenic nerve (C3-C5, innervates the diaphragm) or the vagus nerve (CN X). Acute (<48 hours): most are benign and self-limiting, triggered by: carbonated beverages, eating too quickly, spicy foods, alcohol, sudden temperature change (hot then cold), excitement, stress. Persistent (>48 hours) or intractable (>1 month): underlying pathology should be sought. Causes: (1) Gastrointestinal: GORD, hiatus hernia, gastric distension, gastritis, hepatitis, pancreatitis, bowel obstruction. (2) Central nervous system: stroke, multiple sclerosis, brainstem tumor, encephalitis, meningitis. (3) Thoracic: pneumonia, pleurisy, pericarditis, myocardial infarction, aortic aneurysm, mediastinal mass, lung cancer. (4) Metabolic: uraemia, hyponatraemia, diabetes, alcohol use disorder. (5) Drugs: corticosteroids, benzodiazepines, barbiturates, methyldopa, chemotherapy (platinum-based). (6) Psychogenic: anxiety, conversion disorder, malingering. Management: acute: breath-holding, drinking cold water, Valsalva manoeuvre, pulling knees to chest, eating a spoonful of sugar, or using a paper bag (rebreathing). Persistent/intractable: chlorpromazine (dopamine antagonist, 25-50 mg IM/IV or 25 mg PO TID, first-line, FDA approved), metoclopramide, baclofen, gabapentin, nifedipine, omeprazole. In severe cases: phrenic nerve block or surgical phrenic nerve crush (rare).

\medterm{Hiccups (Singultus)}-- Involuntary, spasmodic contractions of the diaphragm and intercostal muscles followed by sudden glottic closure, producing the characteristic "hic" sound. The reflex arc involves the phrenic nerve (C3-C5), vagus nerve, and brainstem centers. Acute hiccups (<48 hours) are common and benign, triggered by gastric distension (overeating, carbonated beverages, aerophagia), temperature changes (hot/cold food or drink), alcohol, excitement, or stress. Persistent hiccups (>48 hours) and intractable hiccups (>1 month) require evaluation for underlying pathology: gastrointestinal (GERD, hiatal hernia, gastritis), neurologic (stroke, brainstem tumor, multiple sclerosis, meningitis), thoracic (pneumonia, pleurisy, pericarditis, myocardial infarction), metabolic (uremia, hyponatremia, diabetes), or drug-induced (dexamethasone, benzodiazepines, opioids, barbiturates, chemotherapy). Imaging: chest X-ray, CT head/chest/abdomen, upper endoscopy, and brain MRI when neurologic signs are present. Treatment for acute hiccups: vagal maneuvers (breath-holding, Valsalva, drinking water rapidly, swallowing dry sugar, pulling knees to chest, carotid sinus massage), and rebreathing into a bag (raises CO2, may suppress the reflex arc). Pharmacotherapy for persistent hiccups: chlorpromazine (first-line, FDA-approved), metoclopramide, baclofen, gabapentin, or haloperidol; severe refractory cases may require phrenic nerve block or vagus nerve stimulation.

\medterm{Internal Hernia}-- Herniation of bowel through a mesenteric defect common after Roux-en-Y gastric bypass creating a defect at the jejunojejunostomy. May cause intermittent or acute small bowel obstruction. Presents with episodic abdominal pain, nausea, and vomiting. Diagnosed by CT showing clustered dilated small bowel loops and the swirling pattern of mesenteric vessels. Treated surgically with reduction of the hernia and closure of the mesenteric defect. Internal hernias after bariatric surgery are a surgical emergency requiring prompt operative intervention to prevent bowel ischemia and infarction.

\medterm{Intestinal Perforation}-- Full-thickness breach in the bowel wall allowing luminal contents to enter the peritoneal cavity. Causes include diverticulitis, appendicitis, peptic ulcer disease, and trauma. Presents with acute abdomen, peritonitis, and free air on imaging. A surgical emergency requiring emergent laparotomy with primary repair or resection depending on location and degree of contamination. Postoperative management includes broad-spectrum antibiotics and intensive care support.

\medterm{Intestinal Lymphoma}-- Malignant proliferation of lymphoid cells arising in the small or large intestine. Important types include enteropathy-associated T-cell lymphoma, monomorphic epitheliotropic intestinal T-cell lymphoma, and diffuse large B-cell lymphoma. Presentations include abdominal pain, diarrhoea, weight loss, gastrointestinal bleeding, obstruction, perforation, or an incidental mass. Diagnosis requires endoscopic or surgical biopsy with histology, immunophenotyping, and staging imaging.

\medterm{Intramural Hematoma}-- Collection of blood within the wall of the gastrointestinal tract most commonly the duodenum. Caused by trauma, anticoagulation, or coagulopathy. Presents with abdominal pain, vomiting, and sometimes obstructive symptoms. Diagnosis is by CT showing circumferential wall thickening with high-attenuation hematoma. Managed conservatively with bowel rest, nasogastric decompression, and correction of coagulopathy. Most resolve spontaneously without surgical intervention.

\medterm{Lynch Syndrome}-- Hereditary nonpolyposis colorectal cancer syndrome caused by mismatch repair gene mutations including MLH1, MSH2, MSH6, and PMS2 with autosomal dominant inheritance. It increases the risk of colorectal, endometrial, ovarian, gastric, and urinary tract cancers. Diagnosis is suggested by microsatellite instability testing and immunohistochemistry showing loss of mismatch repair proteins, confirmed by genetic testing. Management involves intensive colonoscopic surveillance every one to two years starting at age 20-25, and prophylactic hysterectomy and salpingo-oophorectomy after completion of childbearing for women with Lynch syndrome.

\medterm{Mallory-Weiss Tear}-- Longitudinal mucosal laceration at the gastroesophageal junction caused by forceful vomiting or retching. It typically presents with hematemesis. Most bleeds spontaneously within twenty-four to forty-eight hours. Diagnosis is made by upper endoscopy showing a linear mucosal tear. Management includes observation, proton pump inhibitors, and endoscopic hemostasis with clips or thermal therapy if active bleeding is identified. Recurrence is uncommon and the prognosis is generally excellent.

\medterm{Meckel Diverticulum}-- Congenital remnant of the omphalomesenteric duct located on the antimesenteric border of the distal ileum approximately sixty centimeters from the ileocecal valve. It follows the rule of twos: occurs in approximately two percent of the population, located approximately two feet from the ileocecal valve, approximately two inches in length, and may contain two types of ectopic tissue most commonly gastric and pancreatic mucosa. The most common presentation in children is painless lower gastrointestinal bleeding from ulceration of ectopic gastric mucosa.

\medterm{Meconium Peritonitis}-- Sterile chemical peritonitis resulting from in utero intestinal perforation allowing meconium to spill into the peritoneal cavity. Most commonly associated with cystic fibrosis causing inspissated meconium obstruction, intestinal atresia, and volvulus. May be detected prenatally as fetal ascites with echogenic peritoneal calcifications. After birth presents with abdominal distension, vomiting, and signs of peritonitis. Diagnosis is by abdominal radiography showing intraperitoneal calcifications. Treatment is supportive with bowel rest, nasogastric decompression, and parenteral nutrition; surgery is reserved for ongoing perforation or obstruction.

\medterm{melena}-- Passage of black, tarry, foul-smelling stools due to the presence of digested blood, indicating upper gastrointestinal bleeding (usually proximal to the ligament of Treitz). The black color results from oxidation of haemoglobin to haematin by intestinal enzymes and bacteria. At least 200 mL of blood is typically required in the upper GI tract to produce melena. Common causes: peptic ulcer disease (most common), esophageal varices, Mallory-Weiss tear, erosive gastritis, Dieulafoy lesion, and gastric cancer. melena may persist for several days after a single bleeding episode. Differentiation from black stools caused by iron supplements, bismuth, or liquorice is important (these are not tarry and test negative for occult blood). Management: same as for hematemesis (resuscitation, PPI, urgent upper endoscopy).

\medterm{Nonalcoholic Fatty Liver Disease}-- NAFLD: hepatic steatosis (>5 percent of hepatocytes) in the absence of significant alcohol use, associated with the metabolic syndrome (obesity, type 2 diabetes, dyslipidemia, hypertension). Spectrum: simple steatosis (non-progressive in most) to nonalcoholic steatohepatitis (NASH—steatosis + inflammation + hepatocyte injury) leading to fibrosis, cirrhosis, and hepatocellular carcinoma. Often asymptomatic; elevated ALT/AST (mild) or discovered incidentally. Diagnosis: imaging (ultrasound—fatty liver), elastography (fibrosis), liver biopsy (NASH diagnosis). Management: weight loss (7-10 percent), exercise, control of metabolic comorbidities; pioglitazone/vitamin E in selected NASH; avoid alcohol. Now termed metabolic dysfunction-associated steatotic liver disease (MASLD).

\medterm{Occult Blood}-- Blood present in feces that is not visible to the naked eye, detectable only by chemical or immunochemical testing. fecal occult blood testing (FOBT) is used for colorectal cancer screening. Guaiac-based FOBT (gFOBT) detects haem peroxidase activity (requires dietary restriction—avoid red meat, horseradish, vitamin C). fecal immunochemical test (FIT) uses antibodies specific for human haemoglobin (no dietary restriction, more sensitive, quantitative; threshold typically 10--20 \(\mu\)g Hb/g feces). Positive FIT requires follow-up colonoscopy. Screening with FIT every 1--2 years reduces colorectal cancer mortality by 15--33\%. False positives can occur with upper GI bleeding, NSAIDs, anticoagulants, and menstrual blood. Other causes of occult GI bleeding include peptic ulcer disease, angiodysplasia, and inflammatory bowel disease.

\medterm{Occult Gastrointestinal Bleeding}-- Chronic or intermittent gastrointestinal blood loss that is not visible to the patient and may not produce melaena or haematochezia. It commonly presents with iron-deficiency anaemia or a positive faecal blood test. Causes include peptic ulcer disease, gastrointestinal malignancy, angiodysplasia, inflammatory bowel disease, and small-bowel lesions. Evaluation may include blood counts and iron studies, upper and lower endoscopy, and small-bowel assessment when initial investigations are unrevealing.

\medterm{Odynophagia}-- Painful swallowing, characterized by a burning, squeezing, or sticking sensation behind the sternum during deglutition. Strongly suggests mucosal inflammation or ulceration of the esophagus. Major causes: Infectious Esophagitis (Candida albicans, Herpes Simplex Virus, Cytomegalovirus in immunocompromised/HIV patients), Pill-Induced Esophagitis (tetracyclines, bisphosphonates, NSAIDs, iron), Radiation Esophagitis, and severe GERD. Evaluation: upper endoscopy (EGD) with biopsy.

\medterm{Oesophageal Ulcer}-- A mucosal ulceration of the oesophagus caused by gastro-oesophageal reflux, infection, medication or corrosive injury, inflammatory disease, or malignancy. It may cause odynophagia, dysphagia, chest pain, haematemesis, or melaena. Diagnosis is by upper endoscopy with biopsy when indicated; complications include bleeding, perforation, and stricture formation.

\medterm{Oesophagitis}-- Inflammation of the oesophageal mucosa. Common causes include gastro-oesophageal reflux, infection in immunocompromised patients, medication-related injury, eosinophilic inflammation, radiation, and caustic ingestion. Symptoms include heartburn, odynophagia, dysphagia, and chest pain. Diagnosis is usually by upper endoscopy, with biopsy when needed to determine the cause.

\medterm{Oesophageal Varices}-- See \hyperlink{Bleeding Esophageal Varices}{bleeding esophageal varices}.

\medterm{Ogilvie Syndrome}-- Acute colonic pseudo-obstruction characterized by massive colonic dilation in the absence of mechanical obstruction typically occurring in hospitalized patients with serious comorbidities including post-surgical states, severe infections, and medication effects. Pathophysiology involves autonomic imbalance with sympathetic overactivity causing inhibition of colonic motility. The cecum is disproportionately dilated carrying risk of perforation which occurs when the cecal diameter exceeds twelve centimetres, with management including nasogastric decompression, correction of electrolyte abnormalities, discontinuation of exacerbating medications, neostigmine administration, and colonoscopic decompression for refractory cases.

\medterm{Peptic Ulcer Disease}-- Condition characterized by ulceration of the gastrointestinal mucosa extending through the muscularis mucosae into the submucosa or deeper layers, occurring most commonly in the duodenal bulb and gastric antrum. The two predominant etiologies are Helicobacter pylori infection and nonsteroidal anti-inflammatory drug use, which together account for the vast majority of peptic ulcers. Duodenal ulcers are strongly associated with H. pylori and present with epigastric pain that is classically relieved by food and antacids, whereas gastric ulcers cause pain that is worsened by eating.

\medterm{Peptic Ulcer Bleeding}-- Haemorrhage caused by erosion of an artery or other blood vessel beneath a gastric or duodenal peptic ulcer. It is a common cause of upper gastrointestinal bleeding and may present with haematemesis, coffee-ground vomiting, melaena, or haemodynamic instability. Diagnosis and risk assessment require urgent clinical evaluation and upper endoscopy; complications include shock, recurrent bleeding, and death.

\medterm{Peutz-Jeghers Syndrome}-- An autosomal dominant neurocutaneous disorder caused by germline loss-of-function mutations in the STK11 (LKB1) tumor suppressor gene on chromosome 19p13. Characterized by: 1) Mucocutaneous hyperpigmented macules (dark brown/blue freckle-like spots on lips, buccal mucosa, perioral skin, and digits), and 2) Multiple hamartomatous polyps throughout the gastrointestinal tract (most dense in small intestine/jejunum). Complications: recurrent intussusception, GI bleeding, and markedly increased lifetime risk of GI (colorectal, gastric, pancreatic) and extra-GI malignancies (breast, ovarian, uterine). Surveillance: capsule endoscopy, colonoscopy, mammography, and pelvic MRI.

\medterm{Plummer-Vinson Syndrome}-- Triad of iron deficiency anemia, esophageal web, and dysphagia occurring predominantly in middle-aged women. The esophageal webs are thin mucosal folds typically located in the postcricoid region of the hypopharynx and proximal esophagus. The pathogenesis is thought to involve epithelial atrophy from chronic iron deficiency affecting rapidly dividing mucosal cells. The syndrome is considered a premalignant condition with increased risk of postcricoid squamous cell carcinoma. Diagnosis is established by barium swallow demonstrating the characteristic thin mucosal web in the proximal esophagus, and treatment consists of iron replacement therapy and endoscopic dilation of the web.

\medterm{Portal Hypertension}-- Elevated pressure in the portal venous system defined as hepatic venous pressure gradient exceeding five millimeters of mercury with clinically significant portal hypertension occurring when the gradient exceeds ten. The most common cause is cirrhosis where increased intrahepatic resistance combined with splanchnic vasodilation creates a hyperdynamic circulatory state. Consequences include esophageal variceal hemorrhage, ascites, splenomegaly with hypersplenism, and hepatic encephalopathy. Diagnosis is established by measurement of the hepatic venous pressure gradient, upper gastrointestinal endoscopy to evaluate for varices, and imaging to assess for features of portal hypertension such as splenomegaly and portosystemic collaterals.

\medterm{Portal Hypertensive Gastropathy}-- Diffuse mucosal changes in the stomach due to portal hypertension causing chronic blood loss. Characterized by mosaic pattern and cherry red spots on endoscopy representing dilated mucosal vessels. Most common in patients with advanced cirrhosis and portal hypertension. Managed with beta-blockers to reduce portal pressure and iron supplementation for anemia. Argon plasma coagulation may be used for severe bleeding. The condition is distinct from gastritis and does not respond to acid suppression therapy.

\medterm{Portopulmonary Syndrome}-- Development of pulmonary arterial hypertension in the setting of portal hypertension or liver disease. Results from increased cardiac output and splanchnic vasodilation leading to pulmonary vascular remodeling. Classified as mild, moderate, or based on hemodynamic parameters with mean pulmonary artery pressure greater than twenty millimeters of mercury. Presents with exertional dysphagia, fatigue, and exercise intolerance. Diagnosis requires right heart catheterization confirming elevated pulmonary artery pressure on right heart catheterisation, with treatment including pulmonary vasodilator therapy and liver transplantation in selected cases.

\medterm{Pyloric Stenosis}-- Hypertrophy of the pyloric sphincter muscle causing gastric outlet obstruction most common in young infants. Presents with non-bilious projectile vomiting typically beginning at two to eight weeks of age. Dehydration and metabolic alkalosis with hypochloremia and hypokalemia develop as the condition progresses. Diagnosis is by ultrasound showing pyloric muscle thickness greater than three millimeters and length greater than fifteen millimeters. Treated surgically with Ramstedt pyloromyotomy, which typically resolves with full recovery.

\medterm{Schatzki Ring}-- Thin circumferential mucosal ring located at the squamocolumnar junction of the distal esophagus at the superior margin of a hiatal hernia. The ring is composed of mucosa and submucosa. May be asymptomatic or cause episodic dysphagia particularly to solid food and acute food bolus impaction. Diagnosis is by barium swallow showing a thin circumferential narrowing at the gastroesophageal junction. Upper endoscopy may miss the ring if the esophagus is overdistended with air. Treatment of symptomatic Schatzki rings is endoscopic dilation or disruption of the ring with biopsy forceps or electrocautery.

\medterm{Small Bowel Obstruction}-- Mechanical blockage of the small intestine most commonly caused by adhesions from prior surgery accounting for sixty to seventy-five percent, followed by hernias and malignancy. Clinical presentation includes colicky abdominal pain, bilious vomiting, distension, and obstipation. Diagnosis is established by abdominal CT demonstrating the transition point. Management involves nasogastric decompression, intravenous fluids, and monitoring for ischemia. Complete obstruction with strangulation require emergency laparotomy with resection of non-viable bowel.

\medterm{Soy Protein Intolerance}-- Adverse gastrointestinal reaction to soy proteins, most often seen in infants and young children. It may cause vomiting, diarrhoea, abdominal pain, blood or mucus in the stool, feeding difficulty, or poor growth. Diagnosis is clinical, based on the history and improvement after supervised avoidance, with recurrence on reintroduction supporting the diagnosis. It should be distinguished from immediate IgE-mediated soy allergy and from other causes of chronic gastrointestinal symptoms.

\medterm{Somatostatinoma}-- Rare pancreatic or duodenal tumor producing somatostatin causing diabetes, gallstones, steatorrhea, and hypochlorhydria from suppression of insulin, gallbladder contraction, pancreatic enzymes, and gastric acid respectively. Diagnosed by elevated somatostatin levels. Treatment is surgical resection for localized disease. Somatostatin analogs may control symptoms but do not treat the underlying tumor. The prognosis depends on tumor stage and completeness of surgical resection.

\medterm{Steatorrhea}-- The passage of excess fat in the stool, causing greasy, pale, foul-smelling, floating stools—a hallmark of fat malabsorption. Causes: pancreatic insufficiency (chronic pancreatitis, cystic fibrosis—impaired lipase), biliary obstruction (gallstones, cholangiocarcinoma—reduced bile salts), mucosal disease (celiac disease, tropical sprue), and lymphatic obstruction. Associated symptoms: weight loss, malnutrition, fat-soluble vitamin (A, D, E, K) deficiencies. Diagnosis: fecal fat (qualitative or 72-hour quantitative), elastase (pancreatic), anti-transglutaminase (celiac), imaging; management treats the underlying cause with pancreatic enzyme replacement and diet.

\medterm{Stomach Cancer}-- Gastric adenocarcinoma, a malignancy of the stomach; a leading cause of cancer death worldwide (more common in East Asia). Risk factors: Helicobacter pylori infection (major), diet (smoked/preserved foods, nitrates), smoking, pernicious anemia, familial syndromes (HDGC, Lynch). Types: intestinal (distal, H. pylori-associated, from chronic gastritis/intestinal metaplasia) and diffuse (limits plastica). Features: dyspepsia, early satiety, weight loss, anemia, melena, vomiting; often advanced at diagnosis. Diagnosis: endoscopy with biopsy, staging CT/endoscopic ultrasound. Management: surgical resection (± neoadjuvant chemo), chemotherapy, targeted therapy (HER2, PD-1); prognosis depends on stage.

\medterm{Stomach Ulcer}-- Peptic ulcer in the gastric mucosa associated with H. See \hyperlink{Gastric Ulcer}{Gastric Ulcer}.

\medterm{Variceal Hemorrhage}-- Life-threatening complication of portal hypertension involving rupture of dilated esophageal or gastric varices. The risk is determined by variceal size, red color signs, liver disease severity, and portal hypertension degree. Acute management involves hemodynamic resuscitation, vasoactive drugs, prophylactic antibiotics, and emergent endoscopic band ligation. Transjugular intrahepatic portosystemic shunt is indicated for refractory bleeding. Prevention involves screening endoscopy and prophylactic non-selective beta-blockers or endoscopic variceal ligation for primary prophylaxis, depending on variceal size and patient tolerance.

\medterm{Whipple Procedure}-- Pancreaticoduodenectomy for cancers of the pancreatic head, distal bile duct, or ampulla. Involves removal of the pancreatic head, duodenum, distal bile duct, and gallbladder with three anastomoses: hepaticojejunostomy, pancreaticojejunostomy, and gastrojejunostomy. The pylorus-preserving variant maintains the entire stomach. Mortality at high-volume centers is approximately two to five percent. Complications include pancreatic fistula, delayed gastric emptying, and hemorrhage. Adjuvant chemotherapy with gemcitabine or 5-fluorouracil-based regimens is given after resection to improve survival, but the prognosis remains poor.

\medterm{Wireless Motility Capsule}-- Ingestible capsule that measures pH, pressure, and temperature throughout the gastrointestinal tract. Provides regional transit times through the stomach, small bowel, and colon in a single test. The pH sensor identifies transition points between segments based on characteristic pH changes. Useful for evaluating gastric emptying, small bowel transit, and colonic transit in patients with suspected motility disorders. Non-invasive alternative to scintigraphy with the advantage of providing whole gut transit assessment in a single non-invasive study.

\medterm{Zenker Diverticulum}-- Pulsion diverticulum at the junction of the thyropharyngeus and cricopharyngeus muscles in the proximal esophagus. It presents with dysphagia, regurgitation of undigested food, halitosis, and aspiration. Diagnosis is by barium swallow showing the characteristic outpouching. Treatment is surgical with diverticulectomy or cricopharyngeal myotomy. The endoscopic approach with cricopharyngeal myotomy is increasingly used. Complications include aspiration pneumonia and, rarely, perforation.

\medterm{Zollinger-Ellison Syndrome}-- Gastrin-secreting tumor called gastrinoma causing gastric acid hypersecretion, severe peptic ulcers, and diarrhea. Most tumors are located in the pancreas or duodenum with approximately sixty percent being malignant. Presents with multiple or refractory peptic ulcers, diarrhea, and esophageal reflux. Diagnosed by elevated fasting gastrin levels greater than one thousand picograms per milliliter and positive secretin stimulation test. Treatment involves high-dose proton pump inhibitors and surgical resection for localized gastrinoma, with a cure rate of approximately 30 percent.

\medterm{Abdominal Compartment Syndrome}-- Condition defined by sustained intra-abdominal pressure greater than twenty millimeters of mercury associated with new organ dysfunction. Causes include massive ascites, bowel distension from obstruction or ileus, hemorrhage, and post-surgical conditions particularly after large abdominal wall hernia repairs. Elevated intra-abdominal pressure impairs venous return reducing cardiac output, restricts diaphragmatic excursion causing respiratory failure, compresses the renal veins causing oliguria, and can progress to abdominal compartment syndrome with multi-organ failure.

\medterm{Abdominal Pain}-- Discomfort within the abdomen classified as visceral (dull, midline, poorly localized), parietal (sharp, well localized, worsened by movement), or referred from a distant source. Characterization by SOCRATES guides the differential, which varies by location---right lower quadrant suggests appendicitis; epigastric, pancreatitis or peptic ulcer. Distinguishing surgical from medical causes relies on history, examination, and targeted imaging and laboratory tests.

\medterm{Acute Liver Failure}-- Rapid deterioration of liver function in a previously healthy individual without prior liver disease. Causes include acetaminophen toxicity, viral hepatitis, drug reactions, autoimmune hepatitis, and Wilson disease. Features include coagulopathy, encephalopathy, and multiorgan failure. Presents with jaundice, confusion, and bleeding. Management in intensive care includes monitoring, treatment of reversible causes, and consideration for liver transplantation which is the definitive treatment for irreversible acute liver failure.

\medterm{Adenoma-Carcinoma Sequence}-- Established model of stepwise colorectal carcinogenesis progressing from normal epithelium through adenomas to carcinoma over ten to fifteen years. The sequence begins with APC tumor suppressor gene inactivation, followed by KRAS oncogene activation promoting adenoma growth, and TP53 inactivation associated with progression to invasive carcinoma. This model explains why polyp removal during colonoscopy prevents colorectal cancer. It accounts for seventy to eighty percent of colorectal cancers. The pathway underscores the importance of screening colonoscopy and polypectomy for cancer prevention.

\medterm{Anal Abscess}-- Collection of pus in the perianal or ischiorectal space usually from infected anal glands. Presents with throbbing anal pain, swelling, fever, and difficulty sitting. Diagnosed clinically with tenderness and fluctuance. Treatment is incision and drainage under local anesthesia. Antibiotics are not routinely required for simple abscesses. Anoscopy should be performed to evaluate for an underlying fistula. Recurrence is common if an associated fistula is not identified and treated.

\medterm{Anal Fissure}-- Linear tear in the distal anal mucosa usually posterior in the midline. Presents with painful defecation and bright red blood on wiping or in the toilet. Associated with constipation, hard stools, and chronic diarrhea. The internal anal sphincter spasm perpetuates the fissure by reducing blood flow to the posterior midline. Managed with stool softeners, dietary fiber, topical nitrates or calcium channel blockers, botulinum toxin injection, and lateral internal sphincterotomy for refractory cases.

\medterm{Anal Fistula}-- Abnormal tract connecting the anal canal to the perianal skin usually following a drained abscess. Presents with recurrent drainage, pain, and pruritus. Parks classification categorizes fistulae by relationship to the sphincter muscles. Diagnosed by examination under anesthesia and MRI. Treated surgically with fistulotomy for low fistulae, seton placement for high fistulae, advancement flap repair, or the LIFT procedure. Crohn-related fistulae require medical optimization with biologic agents before surgical intervention to optimise wound healing and reduce recurrence.

\medterm{Piles (Hemorrhoids)}-- Swollen, inflamed vascular cushions in the anal canal that normally aid continence. Internal hemorrhoids (above the dentate line, painless) present with painless bright red rectal bleeding and prolapse; graded I-IV by degree of prolapse. External hemorrhoids (below the dentate line, covered by anoderm) present with pain, swelling, and thrombosis (acute purple mass). Thrombosed external hemorrhoids cause severe perianal pain. Risk factors: constipation, straining, pregnancy, obesity, aging. Diagnosis: clinical (inspection, digital rectal exam, anoscopy). Treatment: lifestyle modifications (fiber, fluids, avoid straining), topical agents (hydrocortisone, lidocaine), rubber band ligation (I-III), sclerotherapy, infrared coagulation, hemorrhoidectomy (III-IV), or stapled hemorrhoidopexy. Thrombosed external hemorrhoids: excision within 48-72 hours of symptom onset for rapid relief.

\medterm{Anastomotic Leak}-- Breakdown of a surgical anastomosis in the gastrointestinal tract causing spillage of luminal contents into the peritoneal cavity. Presents with fever, tachycardia, abdominal pain, leukocytosis, and peritonitis. Diagnosed by CT with oral or rectal contrast showing extraluminal contrast or fluid collection. Risk factors include malnutrition, steroid use, obesity, and distal anastomosis. Managed with broad-spectrum antibiotics, percutaneous drainage of abscesses, and possible reoperation for diffuse peritonitis.

\medterm{Angiodysplasia}-- Vascular ectasia consisting of dilated fragile blood vessels in the gastrointestinal mucosa and submucosa, most commonly found in the cecum and ascending colon. They are the most common cause of acute lower gastrointestinal bleeding in elderly patients over sixty years of age. Patients typically present with intermittent painless hematochezia or melena and chronic iron deficiency anemia. Diagnosis is best achieved by colonoscopy revealing cherry-red spots. Management of actively bleeding lesions includes endoscopic therapy with cautery, argon plasma coagulation, or clipping.

\medterm{Anorectal Manometry}-- Measurement of pressures in the anal canal and rectum to assess sphincter function and rectal sensation. A pressure-sensing catheter is inserted into the rectum and slowly withdrawn through the anal canal. Measures resting and squeeze pressures reflecting internal and external sphincter function respectively. Rectal sensation is assessed by balloon inflation. Used to evaluate fecal incontinence, chronic constipation, and to screen for Hirschsprung disease in children. The study helps guide treatment decisions for these conditions.

\medterm{Anorexia}-- Loss of appetite, distinct from anorexia nervosa (the eating disorder). Causes include: infections (most common), malignancies (especially GI, pancreatic, lung), chronic diseases (CKD, COPD, heart failure, liver disease), medications (chemotherapy, opioids, antibiotics), psychological disorders (depression, anxiety), and endocrine disorders (Addison's disease, hypothyroidism). Pathophysiology involves cytokines (TNF-$\alpha$, IL-6), leptin dysregulation, and altered hypothalamic appetite centers. Consequences: malnutrition, weight loss, sarcopenia, and immune dysfunction. Management: treat underlying cause, nutritional support (oral supplements, enteral feeding if severe), appetite stimulants (megestrol acetate, corticosteroids, dronabinol), and dietary modifications (small frequent meals, preferred foods).

\medterm{APACHE II Score}-- Severity scoring system for critically ill patients incorporating physiological variables, age, and chronic health status to generate scores from zero to seventy-one. Higher scores indicate greater severity and higher mortality. In acute pancreatitis, scores of eight or greater indicate severe disease. The score is calculated within twenty-four hours of ICU admission. It is validated across multiple ICU populations and used for severity assessment, quality improvement, and research. Limitations include the complexity of calculation and the need for accurate physiological data.

\medterm{Appendicitis}-- Inflammation of the vermiform appendix, most commonly due to luminal obstruction by fecalith, lymphoid hyperplasia, or foreign body. Presents with periumbilical pain migrating to the right lower quadrant, nausea, anorexia, and low-grade fever. Diagnosis: clinical, elevated WBC, CT scan. Treatment: laparoscopic appendectomy. Complications: perforation, peritonitis, abscess formation if untreated.

\medterm{Ascites}-- Pathological accumulation of fluid within the peritoneal cavity representing the most common decompensating event in cirrhosis. The pathogenesis involves splanchnic arterial vasodilation reducing effective arterial blood volume and activating neurohormonal mechanisms leading to sodium and water retention. Diagnosis is established by clinical examination and confirmed by ultrasonography. Diagnostic paracentesis determines the serum-ascites albumin gradient, with a gradient of one point one grams per decilitre or greater indicates portal hypertension.

\medterm{Autoimmune Hepatitis}-- Chronic hepatitis caused by immune-mediated destruction of hepatocytes characterized by interface hepatitis with lymphoplasmacytic infiltrate, elevated transaminases, hypergammaglobulinemia, and autoantibodies. Type 1 autoimmune hepatitis is associated with anti-nuclear and anti-smooth muscle antibodies while type 2 is associated with anti-liver kidney microsomal antibodies. Diagnosis is based on the International Autoimmune Hepatitis Group scoring system. Treatment involves corticosteroids with azathioprine as steroid-sparing agents for maintenance therapy.

\medterm{Autoimmune Pancreatitis}-- Chronic inflammatory condition characterized by dense lymphoplasmacytic infiltrate with IgG4-positive plasma cells representing the pancreatic manifestation of IgG4-related disease. Type 1 is more common affecting older men with elevated serum IgG4 and extrapancreatic involvement. Type 2 affects younger patients with higher prevalence of inflammatory bowel disease. Clinical presentation includes painless obstructive jaundice and new-onset diabetes. Imaging shows a sausage-shaped pancreas with peripancreatic fat stranding and capsule-like rim enhancement.

\medterm{Bile Salt Malabsorption}-- Condition characterized by impaired reabsorption of bile acids in the terminal ileum leading to excessive bile acids entering the colon where they stimulate water and electrolyte secretion causing chronic watery diarrhea. Normally approximately ninety-five percent of bile acids are reabsorbed in the terminal ileum through the apical sodium-dependent bile acid transporter. When this mechanism is disrupted by surgical resection, inflammatory bowel disease, or idiopathic causes, excess bile acids in the colon stimulate colonic secretion causing watery diarrhea, which responds to bile acid sequestrants such as colestyramine.

\medterm{Biliary Atresia}-- A progressive, obliterative cholangiopathy of the extrahepatic bile ducts in neonates, leading to biliary obstruction, cholestasis, and ultimately biliary cirrhosis if untreated. It is the most common cause of neonatal jaundice requiring surgery and the leading indication for pediatric liver transplantation, with an incidence of approximately 1 in 10,000 to 15,000 live births. The etiology is unknown but likely involves an inflammatory insult (possible viral trigger: rotavirus, reovirus) superimposed on a developmental susceptibility of the neonatal bile duct epithelium.

\medterm{Biliary Dyskinesia}-- Functional gallbladder disorder with impaired gallbladder ejection fraction in the absence of gallstones. Presents with biliary colic-like right upper quadrant pain. Diagnosed by CCK-HIDA scan showing ejection fraction less than thirty-five percent. The diagnosis requires exclusion of other causes of biliary-type pain. Treatment is cholecystectomy, which provides symptom relief in approximately eighty percent of patients with confirmed biliary dyskinesia.

\medterm{Bowel Incontinence}-- The involuntary loss of solid or liquid stool through the anus, also known as fecal incontinence, affecting approximately 2-15 percent of the general population and increasing in prevalence with age. Causes are classified by mechanism: structural sphincter damage (obstetric anal sphincter injury from vaginal delivery, the most common cause in women; anorectal surgery; trauma), neurogenic dysfunction (diabetic neuropathy, spinal cord injury, stroke, multiple sclerosis, cauda equina syndrome), and cognitive impairment.

\medterm{Budd-Chiari Syndrome}-- Obstruction of hepatic venous outflow at any level from small hepatic venules to the inferior vena cava junction with the right atrium. The most common predisposing condition is myeloproliferative disorders accounting for forty to fifty percent of cases. Other risk factors include thrombophilia and oral contraceptive use. Clinical presentation varies from acute right upper quadrant pain with ascites to chronic presentation mimicking cirrhosis. Diagnosis is established by Doppler ultrasonography showing absent or reversed hepatic venous flow, and confirmed by venography.

\medterm{C282Y Mutation}-- Most common genetic cause of hereditary hemochromatosis, a C to T transition in the HFE gene resulting in cysteine to tyrosine substitution at position 282. Homozygosity occurs in approximately one in two hundred individuals of Northern European descent. The mutation disrupts HFE protein interaction with transferrin receptor leading to inappropriately low hepcidin production and unregulated iron absorption. Penetrance is incomplete with only approximately twenty-eight percent of male and one per cent of female C282Y homozygotes develop clinical iron overload.

\medterm{Cambridge Classification}-- ERCP-based grading system for chronic pancreatitis severity evaluating the main pancreatic duct, branches, and common bile duct. Normal pancreas shows smooth duct without abnormalities. Equivocal shows fewer than three abnormal side branches. Mild shows three or more abnormal side branches. Moderate shows abnormal main duct. Severe shows calculi, cavities, or strictures. The classification has been largely superseded by MRCP and endoscopic ultrasound which provide both ductal and parenchymal assessment.

\medterm{Capsule Endoscopy}-- Non-invasive diagnostic technique using a swallowable wireless endoscopic device to capture images of the small intestinal mucosa. The capsule contains a miniature video camera and transmitter that sends images to a receiver worn on the patient's belt. The procedure is primarily used to evaluate the small bowel for obscure gastrointestinal bleeding, suspected Crohn disease, small bowel tumors, and celiac disease. The patient ingests the capsule after an overnight fast and it is passed naturally in the stool within 24 to 48 hours.

\medterm{Carcinoid Syndrome}-- Serotonin-secreting neuroendocrine tumor causing flushing, diarrhea, bronchospasm, and right-sided heart disease. Occurs when vasoactive substances bypass hepatic metabolism through venous drainage or liver metastases. Diagnosis is with twenty-four hour urine 5-hydroxyindoleacetic acid or serum chromogranin A. Octreotide controls symptoms by suppressing hormone release. Surgical debulking reduces tumor burden. Telotristat ethyl inhibits serotonin synthesis for refractory diarrhea. Right-sided valvular heart disease, particularly tricuspid regurgitation and pulmonary stenosis, is managed with valve replacement in severe cases.

\medterm{Celiac Crisis}-- Acute life-threatening presentation of celiac disease characterized by severe dehydration, electrolyte imbalances, metabolic acidosis, hypocalcemic tetany, and hypovolemic shock. Most commonly occurs in infants and young children after a viral illness or introduction of gluten-containing foods. Presents with profuse watery diarrhea, vomiting, abdominal distension, and signs of severe dehydration. Treatment requires aggressive fluid and electrolyte resuscitation, correction of hypocalcemia and hypoglycaemia, and administration of intravenous corticosteroids.

\medterm{Celiac Disease}-- An immune-mediated enteropathy triggered by gluten ingestion in genetically predisposed individuals (HLA-DQ2/DQ8). Gliadin peptides from wheat, barley, and rye activate an autoimmune response against tissue transglutaminase (tTG), causing villous atrophy, crypt hyperplasia, and intraepithelial lymphocytosis in the small intestine. Presents with diarrhea, steatorrhea, weight loss, bloating, iron deficiency anemia, and failure to thrive in children. Extraintestinal: dermatitis herpetiformis, osteoporosis, infertility, and neuropathy. Diagnosis: positive tTG-IgA, EMA, and duodenal biopsy showing villous atrophy. Management: strict lifelong gluten-free diet. Complications: refractory celiac disease, enteropathy-associated T-cell lymphoma (rare).

\medterm{Celiac Trunk Stenosis}-- Narrowing of the celiac trunk typically caused by extrinsic compression from the median arcuate ligament of the diaphragm, known as median arcuate ligament syndrome or celiac artery compression syndrome. Causes chronic abdominal pain particularly after eating and weight loss due to postprandial mesenteric ischemia. Diagnosis is by CT angiography or conventional angiography showing characteristic hook-shaped narrowing of the celiac trunk and collateral vessel formation. Treatment is surgical division of the median arcuate ligament with or without celiac artery reconstruction.

\medterm{Charcot Triad}-- Clinical constellation of fever, right upper quadrant pain, and jaundice representing the classic presentation of ascending cholangitis, infection of the biliary tree from common bile duct obstruction by a gallstone. The full triad is found in approximately fifty to seventy percent of patients. When hypotension and altered mental status are added, it constitutes Reynolds pentad indicating suppurative cholangitis. Initial management involves intravenous fluids, broad-spectrum antibiotics, and urgent biliary drainage by ERCP with endoscopic sphincterotomy.

\medterm{Charcot Triad (Cholangitis)}-- Classic diagnostic triad of acute ascending cholangitis: (1) Fever, (2) Right Upper Quadrant abdominal pain, and (3) Jaundice. Reynolds Pentad adds Hypotension and Altered Mental Status. \textit{See Acute Cholangitis.}

\medterm{Child-Pugh Classification}-- Clinical scoring system assessing prognosis of chronic liver disease by evaluating five parameters: total bilirubin, serum albumin, international normalized ratio, ascites, and hepatic encephalopathy. Each parameter is scored one to three points. Child-Pugh A with five to six points represents well-compensated cirrhosis with one-year survival exceeding ninety-five percent. Child-Pugh B with seven to nine points represents significant functional compromise with one-year survival approximately eighty percent. Child-Pugh C with ten to fifteen points represents decompensated cirrhosis with one-year survival of approximately fifty percent.

\medterm{Cholangitis}-- Bacterial infection of the biliary tree, most commonly caused by choledocholithiasis (common bile duct stones) obstructing bile flow, but also from biliary strictures, stent occlusion, and malignancy. The classic clinical presentation is Charcot triad: fever, right upper quadrant pain, and jaundice (present in 50-70 percent of cases). Suppurative cholangitis with Reynolds pentad (Charcot triad plus hypotension and altered mental status) indicates severe disease requiring urgent decompression. Complications include pancreatitis, perforation, bleeding, and cholangitis.

\medterm{Choledochal Cyst}-- Congenital dilatation of the biliary tree classified into five types according to the Todani classification. Type I involving dilation of the common bile duct is the most common. Presents with the classic triad of abdominal pain, jaundice, and palpable mass though all three are present in less than twenty percent of cases. Risk of cholangiocarcinoma is significantly increased. Diagnosed by MRCP or CT showing biliary dilation. Treated with complete surgical excision of the cyst and biliary reconstruction with Roux-en-Y hepaticojejunostomy.

\medterm{Choledocholithiasis}-- The presence of gallstones within the common bile duct, occurring in approximately 10-20 percent of patients with cholelithiasis. Stones may form primarily in the common bile duct (primary stones, brown pigment stones, associated with biliary stasis and infection) or migrate from the gallbladder (secondary stones, most common, identical to gallbladder stones). Clinical presentation ranges from asymptomatic (incidental finding on imaging) to classic biliary colic, obstructive jaundice, and cholangitis. Endoscopic retrograde cholangiopancreatography is both diagnostic and therapeutic, allowing stone extraction after sphincterotomy.

\medterm{Chronic Intestinal Pseudo-Obstruction}-- Chronic condition of recurrent episodes of intestinal obstruction without mechanical blockage due to neuromuscular dysfunction of the gastrointestinal tract. May involve the small bowel, colon, or both. Causes include systemic sclerosis, diabetes, Parkinson disease, amyloidosis, and paraneoplastic syndromes. Presents with abdominal distension, pain, nausea, vomiting, constipation, and malabsorption. Diagnosed by imaging showing dilated bowel loops without transition point, and antroduodenal manometry showing abnormal motility patterns in the small bowel.

\medterm{Chronic Mesenteric Ischemia}-- Recurrent abdominal pain occurring after eating resulting from insufficient blood flow to meet the metabolic demands of the gastrointestinal tract. Most commonly caused by atherosclerotic stenosis of at least two of the three major mesenteric arteries. Classic presentation is sitophobia, fear of eating due to postprandial abdominal pain typically beginning thirty minutes to two hours after meals and lasting one to two hours. Patients develop food avoidance and significant weight loss. Diagnosis is confirmed by CT angiography or conventional angiography demonstrating stenosis of at least two mesenteric vessels.

\medterm{Chronic Pancreatitis}-- Progressive irreversible inflammatory condition characterized by permanent structural damage causing exocrine and endocrine impairment. The most common cause is chronic alcohol consumption. Other causes include hereditary pancreatitis, autoimmune pancreatitis, and metabolic factors. Clinical manifestations include chronic abdominal pain, steatorrhea, and diabetes mellitus. Diagnosis is established by imaging showing calcifications, ductal dilatation, and parenchymal atrophy. The Cambridge classification is used for staging severity.

\medterm{Clostridioides difficile Infection}-- Bacterial infection of the colon caused by C. difficile toxin, usually following antibiotic use that disrupts normal colonic flora. Presents with watery diarrhea, fever, and abdominal pain. Diagnosed by stool toxin assay or nucleic acid amplification test. Treatment involves discontinuation of the offending antibiotic and initiation of oral vancomycin or fidaxomicin. Fecal microbiota transplantation is highly effective for recurrent infection. Complications include toxic megacolon, perforation, and death.

\medterm{CMV Colitis}-- Colonic inflammation caused by cytomegalovirus reactivation in immunocompromised patients including those with HIV and organ transplant recipients. Presents with bloody diarrhea and abdominal pain. Diagnosed by colonoscopy with biopsy showing characteristic owl-eye intranuclear inclusions in enlarged cells. Treated with intravenous ganciclovir or oral valganciclovir. CMV colitis should be considered in any immunocompromised patient with colitis unresponsive to standard therapy.

\medterm{Colitis}-- Inflammation of the colon, with a broad differential diagnosis and diverse clinical presentations. Acute colitis presents with diarrhea (often bloody), abdominal pain, tenesmus, and urgency. The most common causes include infectious colitis (bacterial: Shigella, Salmonella, Campylobacter, E. coli O157:H7, C. difficile; viral: CMV; parasitic: amebiasis), ischemic colitis (from hypoperfusion, typically at watershed areas of splenic flexure and rectosigmoid), and idiopathic inflammatory bowel disease (ulcerative colitis and Crohn disease).

\medterm{Collagenous Colitis}-- Chronic inflammatory bowel disease characterized by thickening of the subepithelial collagen band in the colonic mucosa to greater than ten micrometers normally less than five micrometers with preservation of the epithelial surface. Presents with chronic watery non-bloody diarrhea primarily in middle-aged women. The collagen band disrupts water and electrolyte absorption leading to secretory diarrhea. Diagnosis requires colonoscopy with random colonic biopsies demonstrating the characteristic subepithelial collagen deposition.

\medterm{Colonic Ischemia}-- Decreased blood flow to the colon causing mucosal or transmural injury. Most commonly affects watershed areas particularly the splenic flexure and rectosigmoid junction. Causes include non-occlusive hypoperfusion from cardiogenic shock, vasoconstrictive medications, and vasculitis. Presents with sudden onset crampy left lower quadrant pain followed by bloody diarrhea within twenty-four hours. Most cases are self-limited with mucosal ischemia resolving within one to two weeks. Severe transmural infarction requires surgical resection.

\medterm{Colonoscopy}-- Visualization of the entire colon and terminal ileum using a flexible colonoscope. It is the gold standard for colorectal cancer screening and the preferred method for evaluating colonic pathology. The procedure requires bowel preparation with purgative solution to clear the colon. Therapeutic capabilities include polypectomy, hemostasis, dilation, and stent placement. Complications include perforation, bleeding, and post-polypectomy syndrome. Adequate adenoma detection rate is a key quality metric in colonoscopy.

\medterm{Colorectal Cancer Screening}-- Public health strategy to detect colorectal cancer and precancerous polyps in asymptomatic individuals. Current guidelines recommend screening beginning at age forty-five. Options include fecal immunochemical test annually, multi-target stool DNA test every three years, colonoscopy every ten years, CT colonography every five years, and flexible sigmoidoscopy every five years. The choice depends on patient preference, access, risk assessment, and shared decision-making. Screening reduces cancer incidence and mortality.

\medterm{Constipation}-- A common symptom-based disorder defined by the Rome IV criteria as having two or more of the following for at least 3 months: straining during more than 25 percent of bowel movements, lumpy or hard stools, sensation of incomplete evacuation, sensation of anorectal obstruction, manual maneuvers to facilitate defecation, or fewer than 3 bowel movements per week. Constipation is classified as primary (normal transit, slow transit, or defecatory dysfunction from pelvic floor dyssynergia or structural abnormalities of the anorectum.

\medterm{Courvoisier Sign}-- Clinical finding of a palpable non-tender gallbladder in obstructive jaundice suggesting malignant biliary obstruction. In gallstone disease, the gallbladder is typically fibrotic and not palpable, whereas in malignant obstruction it is distended because the wall is initially normal and compliant. The most common cause is pancreatic head carcinoma. The sign has a positive predictive value of approximately eighty-five to ninety percent. Physical examination requires careful right upper quadrant palpation.

\medterm{Crohn Disease}-- Chronic transmural inflammatory condition that can affect any segment of the gastrointestinal tract from the oral cavity to the perianal region, though the terminal ileum and ileocolonic region are most commonly involved. Unlike ulcerative colitis, Crohn disease is characterized by transmural inflammation that extends through all layers of the bowel wall, segmental or skip lesions with intervening normal bowel, and a cobblestone mucosal appearance resulting from intersecting lines of ulceration. The transmural nature of inflammation predisposes to fistula formation, abscesses, and strictures.

\medterm{Crohn Disease Cobblestoning}-- Endoscopic appearance of the colonic or small bowel mucosa in Crohn disease characterized by a bumpy, irregular surface resembling cobblestones. This distinctive appearance results from intersecting lines of deep linear and serpiginous ulceration that surround areas of relatively preserved or edematous but intact mucosa. The combination of raised mucosal islands between ulcer tracks creates the characteristic cobblestone pattern. This finding is highly suggestive of Crohn disease and reflects the underlying pathology of the disease.

\medterm{CT Colonography}-- CT imaging of the colon after bowel preparation and insufflation with air or carbon dioxide. Used for colorectal cancer screening in patients unable to undergo optical colonoscopy and for evaluating complete colonoscopy when the cecum was not reached. Detects polyps greater than six millimeters with high sensitivity. Advantages include non-invasiveness and examination of extracolonic structures. Disadvantages include radiation exposure, need for adequate bowel preparation, and inability to perform polypectomy or biopsy at the time of imaging.

\medterm{CT Enterography}-- CT imaging of the small bowel after administration of large volumes of neutral oral contrast distending the small bowel with IV contrast. Excellent for evaluating Crohn disease, small bowel tumors, and obscure gastrointestinal bleeding. Demonstrates bowel wall thickening, hyperenhancement, mesenteric fat stranding, lymphadenopathy, and complications including strictures, fistulaes, and abscesses. The technique provides high spatial resolution and rapid acquisition. It is limited by ionizing radiation exposure.

\medterm{Curling Ulcer}-- Stress ulceration of the stomach or duodenum associated with severe burns. Part of the Curling triad with hypovolemia and sepsis. The pathogenesis involves reduced mucosal blood flow during the systemic inflammatory response to burn injury. Prevention with stress ulcer prophylaxis using proton pump inhibitors or H2 blockers is standard in burn unit patients. Treatment of established Curling ulcers involves acid suppression and surgical intervention for complications including hemorrhage and perforation.

\medterm{Defecography}-- Dynamic imaging of the rectum and pelvic floor during defecation using fluoroscopy or MRI. The rectum is filled with contrast paste and the patient performs defecation while images are captured. Evaluates rectal prolapse, intussusception, rectocele, enterocele, and pelvic floor dyssynergia. MR defecography provides superior soft tissue detail without radiation and is the preferred modality when available. The study is particularly useful in evaluating patients with obstructed defecation syndrome, particularly when standard therapies have failed.

\medterm{Diarrhea}-- Frequent passage of loose, watery stools, typically defined as more than 3 bowel movements per day with stool weight exceeding 200 grams. Acute diarrhea (less than 14 days) is most commonly infectious in etiology: viral (norovirus, rotavirus, adenovirus), bacterial (Campylobacter, Salmonella, Shigella, E. coli, C. difficile), and parasitic (Giardia, Cryptosporidium, Entamoeba). Chronic diarrhea (more than 4 weeks) requires evaluation for malabsorption syndromes (celiac disease, pancreatic insufficiency, bile acid malabsorption, and microscopic colitis).

\medterm{Diarrhea, Constipation, Gas, and Bloating}-- These common gastrointestinal symptoms arise from disruptions in normal digestive function caused by diet, infections, medications, or underlying conditions like irritable bowel syndrome. Diarrhea involves increased stool frequency and water content, while constipation involves infrequent or difficult passage of stool. Gas and bloating result from swallowed air or colonic fermentation of undigested food and may cause abdominal distension and discomfort.

\medterm{Diverticular Disease}-- Condition characterized by the presence of diverticula, which are outpouchings of the colonic mucosa and submucosa through defects in the muscularis propria at sites where vasa recta penetrate the bowel wall. Diverticulosis, the presence of diverticula without inflammation, is extremely common in Western populations, affecting more than half of individuals over age sixty, and is thought to result from a combination of low dietary fiber intake, increased intraluminal pressure from straining, and age-related weakening of the colonic wall.

\medterm{Diverticulitis}-- Inflammation and infection of colonic diverticula, a common complication of diverticulosis affecting 10-25 percent of individuals with diverticula. The pathophysiology involves microperforation of a diverticulum with fecal contents entering the pericolic tissues, inciting inflammation and potential abscess formation. Clinical presentation includes acute onset of left lower quadrant abdominal pain (typically constant, localized, worsened by movement), fever, nausea, vomiting, and altered bowel habit, and anorexia.

\medterm{Diverticulosis}-- The presence of diverticula, small mucosal outpouchings through the colonic wall at sites of vascular penetration, most commonly in the sigmoid colon (90 percent of cases). Prevalence increases with age, affecting over 50 percent of individuals older than 60 years and nearly 70 percent by age 85. The pathogenesis involves a combination of low dietary fiber intake leading to increased intraluminal pressure, segmental colonic muscle hypertrophy, and weakening of the colonic wall at points of nutrient arteries at the point of vasa recta penetration through the circular muscle layer.

\medterm{Diverticulosis and Diverticulitis}-- Diverticulosis is a condition in which small, bulging pouches (diverticula) form in the lining of the digestive tract, most commonly in the colon. Diverticulitis occurs when these pouches become inflamed or infected, causing abdominal pain, fever, and changes in bowel habits. Treatment ranges from dietary modification and antibiotics for mild cases to surgical intervention for severe or recurrent episodes.

\medterm{Double-Balloon Enteroscopy}-- Advanced endoscopic technique using an enteroscope equipped with balloons at the tip of both the endoscope and an overtube allowing systematic advancement and visualization of the entire small bowel. The balloons are alternately inflated and deflated to anchor the bowel wall and advance the endoscope. This technique allows deep intubation of the small bowel from both oral and anal approaches. Double-balloon enteroscopy is the most comprehensive method for evaluating small bowel pathology, with the ability to perform biopsies and therapeutic interventions throughout the entire small bowel.

\medterm{Dubin-Johnson syndrome}-- A benign, autosomal recessive disorder of conjugated hyperbilirubinemia caused by mutations in the ABCC2 gene encoding the multidrug resistance-associated protein 2 (MRP2), which transports conjugated bilirubin from hepatocytes into bile. This results in impaired excretion of conjugated bilirubin and other organic anions, leading to elevated serum conjugated bilirubin levels (typically 2-5 mg/dL). The liver has a characteristic gross black appearance due to accumulation of dark pigment (epinephrine metabolites) deposited in the hepatocytes.

\medterm{Dysentery}-- Severe diarrheal illness with blood and mucus in stool, typically caused by Shigella, Salmonella, Campylobacter, or Entamoeba histolytica. Presents with frequent bloody stools, abdominal cramps, tenesmus, and fever. Diagnosis: stool culture, microscopy, PCR. Treatment: rehydration and targeted antibiotics based on pathogen and sensitivity. Complications: toxic megacolon, hemolytic uremic syndrome (Shiga toxin-producing E. coli).

\medterm{Endoscopic Retrograde Cholangiopancreatography}-- Combined endoscopic and fluoroscopic procedure used to visualize and treat disorders of the biliary and pancreatic ducts. A duodenoscope is advanced to the second portion of the duodenum and the ampulla of Vater is cannulated. Contrast dye is injected under fluoroscopic guidance creating cholangiograms and pancreatograms. Therapeutic interventions include endoscopic sphincterotomy, stone extraction using baskets and balloons, stent placement for strictures, and nasobiliary drainage. The most common complication is post-ERCP pancreatitis, occurring in 3-10 percent of cases.

\medterm{Enema}-- The introduction of fluid into the rectum and colon through the anus for diagnostic or therapeutic purposes. Types: cleansing enema (warm water, saline, or soap solution to relieve constipation or prepare for colonoscopy/surgery), retention enema (mineral oil for fecal impaction, corticosteroid or mesalamine for inflammatory bowel disease), barium enema (radiographic contrast study of the colon--now largely replaced by colonoscopy and CT colonography), and medicinal enema (administration of drugs such as laxatives, anti-inflammatories, or anesthetics). Potential complications: electrolyte disturbance (with excessive tap water enemas), bowel perforation (rare, with improper technique), and vagal reaction.

\medterm{Enteritis}-- Inflammation of the small intestine, most commonly caused by infectious agents (viral, bacterial, or parasitic). Viral enteritis (norovirus, rotavirus, adenovirus, astrovirus) presents with acute onset of watery diarrhea, nausea, vomiting, and abdominal cramps, typically self-limited over 2--7 days. Bacterial enteritis (Campylobacter, Salmonella, Shigella, E. coli, Yersinia) more often presents with fever, bloody diarrhea, and tenesmus. Chronic enteritis occurs in Crohn disease, celiac disease, radiation enteritis, and eosinophilic gastroenteritis. Diagnosis: clinical history, stool studies (culture, PCR, ova and parasites), and imaging. Management: rehydration (oral or IV), antiemetics, and antibiotics only for specific bacterial infections.

\medterm{Familial Adenomatous Polyposis}-- Hereditary condition with hundreds to thousands of adenomas throughout the colon caused by germline APC gene mutations with autosomal dominant inheritance. Polyps develop during the second decade and colorectal cancer occurs by age forty if untreated. Diagnosis is by colonoscopy showing numerous polyps and genetic testing. Management requires prophylactic colectomy typically in the second decade. Upper gastrointestinal endoscopy surveillance is needed for duodenal and periampullary polyps. Genetic counselling and testing of at-risk family members are essential components of management.

\medterm{Fatty Liver Disease (MASLD)}-- A condition characterized by the accumulation of excess fat in liver cells, most commonly associated with obesity, insulin resistance, and metabolic syndrome. Non-alcoholic fatty liver disease (NAFLD) is the most common form and ranges from simple steatosis to non-alcoholic steatohepatitis (NASH), which can progress to fibrosis and cirrhosis. Management focuses on weight loss, dietary modification, exercise, and control of metabolic risk factors.

\medterm{Fecal Calprotectin}-- Neutrophil-derived protein measured in stool as a marker of intestinal inflammation. It is stable in stool at room temperature for up to seven days. Useful for differentiating inflammatory from functional bowel disease and monitoring inflammatory bowel disease activity. Levels above fifty micrograms per gram suggest intestinal inflammation. Fecal calprotectin has high sensitivity and specificity for detecting endoscopically active inflammatory bowel disease and can reduce the need for endoscopic surveillance when used as a triage tool.

\medterm{Fecal Elastase}-- Pancreatic enzyme measured in stool as a marker of exocrine pancreatic function. Low levels below two hundred micrograms per gram indicate pancreatic insufficiency, and levels below one hundred are consistent with severe exocrine insufficiency. Useful for diagnosing chronic pancreatitis and cystic fibrosis-related pancreatic insufficiency. The test is not affected by enzyme supplementation. Fecal elastase has good sensitivity and specificity for detecting moderate to severe pancreatic exocrine insufficiency.

\medterm{Fecal Incontinence}-- Inability to control bowel movements ranging from minor soiling to complete loss of continence. Causes include sphincter injury from obstetric trauma or surgery, neuropathy from diabetes or spinal cord injury, diarrhea, rectal prolapse, and cognitive impairment. Diagnosed by anorectal manometry, endoanal ultrasound, and pudendal nerve terminal motor latency testing. Managed with pelvic floor exercises and biofeedback as first-line therapy, dietary modification to optimize stool consistency, bulk-forming laxatives as initial management, with prescription medications such as loperamide for diarrhea-predominant symptoms and linaclotide for constipation-predominant symptoms.

\medterm{Fecal Microbiota Transplantation}-- Transfer of stool from a healthy donor into the gastrointestinal tract of a recipient to restore microbiome diversity. Highly effective for recurrent Clostridioides difficile infection with cure rates exceeding eighty-five percent. Administered via colonoscopy, enema, nasoduodenal tube, or oral capsules. The procedure is generally safe with rare complications including aspiration during administration. Donor screening is essential to prevent transmission of infectious diseases. Research is ongoing into the role of FMT in other conditions such as inflammatory bowel disease and metabolic syndrome.

\medterm{FibroScan}-- Non-invasive transient elastography technique measuring liver stiffness to assess hepatic fibrosis. Values correlate with fibrosis stage with higher values indicating more advanced fibrosis. The technique avoids the need for liver biopsy in many patients. Results may be unreliable in obesity, ascites, and narrow intercostal spaces. Controlled attenuation parameter measures hepatic steatosis simultaneously. FibroScan is useful for screening, diagnosis, and monitoring of liver fibrosis in chronic liver disease, particularly in patients with chronic hepatitis B, non-alcoholic fatty liver disease, and recurrent alcohol-related liver disease.

\medterm{Fistula-in-Ano}-- Abnormal epithelialized tract connecting the anal canal to the perianal skin, representing a common and distressing complication of Crohn disease and also occurring following perianal abscess drainage or from cryptoglandular infection. In Crohn disease, fistulae result from transmural inflammation that penetrates through the bowel wall and tracks to the skin surface. Parks classification categorizes fistulae based on their relationship to the anal sphincter muscles: intersphincteric fistulae course between the internal and external anal sphincters, trans-sphincteric fistulae pass through the external sphincter, supra-sphincteric fistulae loop over the puborectalis, and extrasphincteric fistulae track outside the sphincter complex entirely.

\medterm{Fitz-Hugh-Curtis Syndrome}-- Perihepatitis associated with pelvic inflammatory disease caused by ascension of Neisseria gonorrhoeae or Chlamydia trachomatis from the pelvis to the liver capsule along the right paracolic gutter. Presents with right upper quadrant pain, fever, and tender hepatomegaly often in young sexually active women. May mimic cholecystitis or hepatitis. Diagnosis requires demonstration of pelvic inflammatory disease with cervical cultures or nucleic acid amplification testing and imaging showing perihepatic enhancement on CT or MRI.

\medterm{Focal Nodular Hyperplasia}-- Benign liver lesion with a characteristic central stellate scar visible on imaging. The lesion represents a hyperplastic response to a pre-existing vascular malformation rather than a true neoplasm. Associated with oral contraceptive use though it also occurs in men and children. No risk of hemorrhage or malignant transformation. Diagnosed by imaging showing a well-circumscribed enhancing mass with central scar. Managed expectantly with no treatment required. Differentiation from hepatic adenoma, which carries risk of hemorrhage and malignant transformation and is managed differently.

\medterm{Foodborne Illness}-- Any illness resulting from the ingestion of contaminated food, encompassing infections and intoxications caused by pathogens, toxins, or chemical contaminants. It affects millions annually and poses particular risks to pregnant women, young children, older adults, and immunocompromised individuals. Prevention relies on proper food handling, cooking temperatures, hand hygiene, and safe storage practices.

\medterm{Food Poisoning}-- Illness caused by consuming food contaminated with bacteria, viruses, parasites, or their toxins. Common pathogens include Salmonella, Escherichia coli, Listeria, and norovirus, with symptoms such as nausea, vomiting, diarrhea, abdominal cramps, and fever. Most cases resolve without treatment, but severe infections may require hospitalization and supportive care.

\medterm{Gall-Bladder}-- A pear-shaped sac located beneath the liver on the visceral surface, in the gallbladder fossa between the right and quadrate lobes of the liver. Capacity: approximately 30-50 mL. Functions: (1) concentrates and stores bile (10-fold concentration by absorption of water and electrolytes), (2) releases bile into the duodenum via the cystic duct and common bile duct in response to cholecystokinin (CCK) released after a fatty meal, (3) secretes mucus (protects gallbladder mucosa from bile acids). The gallbladder is absent in some species (rat, horse) and can be removed (cholecystectomy) without significant digestive impairment. Bile composition: bile salts, bilirubin, cholesterol, phospholipids.

\medterm{Gall-Bladder, Diseases of}-- Pathological conditions affecting the gallbladder. Types: (1) Cholelithiasis (gallstones)--most common, affects 10-15\% of adults in Western countries. (2) Acute cholecystitis--inflammation of the gallbladder, usually due to stone obstruction of the cystic duct (90\%). (3) Chronic cholecystitis--persistent inflammation with fibrosis, almost always associated with stones. (4) Acalculous cholecystitis--inflammation without stones, seen in critically ill patients, burns, trauma, total parenteral nutrition. (5) Gallbladder polyps--most are benign cholesterol polyps; adenomatous polyps have malignant potential. (6) Gallbladder carcinoma--rare (1-2\% of gallstone patients), more common with large stones (>3 cm), porcelain gallbladder, and chronic Salmonella typhi carriage. (7) Choledocholithiasis--stones in the common bile duct. (8) Mirizzi syndrome--stone impacted in the cystic duct causing extrinsic compression of the common hepatic duct.

\medterm{Gallstone Ileus}-- Large gallstone eroding through the gallbladder wall into the duodenum causing mechanical small bowel obstruction. Diagnosed by CT showing Rigler triad of small bowel obstruction, pneumobilia, and ectopic gallstone. The most common site of obstruction is the ileum where the bowel caliber narrows. Treatment is surgical with enterolithotomy to relieve the obstruction. Cholecystectomy and fistula repair may be performed simultaneously or as a staged procedure depending on the patient's condition.

\medterm{Gallstones}-- Crystalline deposits that form in the gallbladder, composed predominantly of cholesterol (80 percent of stones in Western populations, yellow-green, radiolucent on plain x-ray, develop when bile is supersaturated with cholesterol) or bilirubin pigment (black pigment stones: sterile bile with high unconjugated bilirubin from hemolysis, cirrhosis, age; brown pigment stones: infected bile from bacterial beta-glucuronidase activity, associated with biliary stasis and parasitic infection in East Asian populations). Also known as cholelithiasis.

\medterm{Gallstones (Cholelithiasis)}-- Crystalline deposits that form in the gallbladder, composed predominantly of cholesterol (80 percent of stones in Western populations, yellow-green, radiolucent on plain x-ray, develop when bile is supersaturated with cholesterol) or bilirubin pigment (black pigment stones: sterile bile with high unconjugated bilirubin from hemolysis, cirrhosis, age; brown pigment stones: infected bile from bacterial beta-glucuronidase activity, associated with biliary stasis and parasitic infection in East Asian populations). Also known as cholelithiasis. See \hyperlink{Gallstones}{Gallstones}.

\medterm{Gastrocnemius}-- The large superficial muscle of the posterior calf, forming the bulk of the calf contour. Origin: medial and lateral femoral condyles. Insertion: calcaneus via the Achilles tendon. Nerve: tibial nerve (S1-S2). Action: plantar flexion of the ankle (standing, walking, running, jumping); also flexes the knee. The gastrocnemius works with the soleus as the triceps surae. Clinical relevance: medial gastrocnemius strain (tennis leg), gastrocnemius rupture, deep vein thrombosis (DVT) may cause calf swelling and must be distinguished from muscle injury. Power: can generate up to 3-4 times body weight during running.

\medterm{Gastroenteritis}-- Inflammation of the stomach and intestinal mucosa, most commonly caused by viral infections (norovirus, rotavirus, adenovirus), bacterial pathogens (E. coli, Salmonella, Campylobacter), or parasites. Presents with diarrhea, vomiting, abdominal cramps, and dehydration. Self-limited in most cases. Management: oral rehydration, antiemetics, and symptomatic care. Antibiotics reserved for specific bacterial infections.

\medterm{Gastrointestinal Stromal Tumor}-- Mesenchymal tumor arising from interstitial cells of Cajal in the gastrointestinal wall most commonly in the stomach. Often driven by gain-of-function mutations in the KIT proto-oncogene or PDGFRA gene. The majority arebut approximately twenty to thirty percent are malignant. Diagnosed by endoscopic ultrasound with fine-needle aspiration and immunohistochemistry showing CD117 or DOG1 positivity. Treated with imatinib a tyrosine kinase inhibitor for unresectable or metastatic disease and surgical resection for localized disease, with imatinib as first-line therapy for advanced or metastatic GIST.

\medterm{Giardiasis}-- Infection with Giardia lamblia (G. intestinalis), a protozoan parasite causing diarrheal illness. Transmission via contaminated water (fecal-oral). Presents with foul-smelling, greasy diarrhea, bloating, abdominal cramps, weight loss, and steatorrhea. Diagnosis: stool antigen testing or microscopy. Treatment: metronidazole or tinidazole. Prevention: water purification, hand hygiene.

\medterm{Gilbert Syndrome}-- A common, benign autosomal recessive hereditary hyperbilirubinemia caused by a promoter mutation in the UGT1A1 gene leading to reduced activity ($\sim 30\%$ of normal) of UDP-glucuronosyltransferase. Characterized by mild, transient unconjugated (indirect) hyperbilirubinemia (total bilirubin typically $<3\text{ mg/dL}$) triggered by physical stress, fasting, illness, dehydration, or strenuous exercise. Asymptomatic with normal liver enzymes (ALT, AST, alkaline phosphatase) and no hemolysis. Excellent prognosis requiring no treatment; reassurance is key.

\medterm{Gluten}-- A complex of storage proteins (prolamins and glutelins) found in wheat (gliadin), barley (hordein), and rye (secalin). Gluten provides elasticity to dough. In susceptible individuals, gluten triggers an autoimmune response (celiac disease) or non-autoimmune sensitivity (non-celiac gluten sensitivity, NCGS). Gluten-free diet is the treatment for celiac disease. Oats contain avenin (structurally different) and are tolerated by most celiac patients if certified gluten-free (avoid cross-contamination). Gluten is also implicated in: dermatitis herpetiformis (cutaneous manifestation of celiac disease), gluten ataxia, and wheat allergy (IgE-mediated, different pathogenesis).

\medterm{Gluten-Sensitive Enteropathy (Celiac Disease)}-- An autoimmune enteropathy triggered by ingestion of gluten (wheat, barley, rye) in genetically predisposed individuals. Incidence: 1\% worldwide, higher in HLA-DQ2 (90\%) and HLA-DQ8 carriers. Pathophysiology: gluten peptides (especially 33-mer alpha-gliadin) are deamidated by tissue transglutaminase (tTG), presented by HLA-DQ2/DQ8, activating T-cell-mediated damage to small intestinal villi (villous atrophy, crypt hyperplasia, increased intraepithelial lymphocytes). Symptoms: (1) Classical: chronic diarrhea, steatorrhea, weight loss, bloating, abdominal pain, failure to thrive (children). (2) Non-classical: iron-deficiency anemia, fatigue, osteoporosis, elevated transaminases, infertility, neuropathy, dermatitis herpetiformis. (3) Asymptomatic: detected by screening in at-risk groups (family history, type 1 diabetes, Down syndrome, autoimmune thyroiditis). Diagnosis: positive serology (IgA anti-tTG, confirm with anti-EMA or anti-DGP), duodenal biopsy (Marsh grade 0-IIIc). Treatment: lifelong strict gluten-free diet. Complications: refractory celiac disease (unresponsive to >12 months of diet), enteropathy-associated T-cell lymphoma (EATL), ulcerative jejunoileitis, increased mortality (hazard ratio 1.2-1.4).

\medterm{Graft-Versus-Host Disease}-- Immunological reaction where donor T cells attack recipient tissues after allogeneic bone marrow transplant. Affects the skin, liver, and gastrointestinal tract. GI manifestations include diarrhea, abdominal pain, nausea, anorexia, and gastrointestinal bleeding. Acute graft-versus-host disease occurs within the first hundred days and chronic occurs after day hundred. Diagnosis is by endoscopy with biopsy showing characteristic apoptosis and crypt abscesses. Treated with corticosteroids, calcineurin inhibitors (tacrolimus), mycophenolate mofetil, and anti-TNF agents.

\medterm{Griping (see Colic)}-- Severe abdominal pain, typically spasmodic in nature. See \hyperlink{Colic}{Colic}.

\medterm{Gut Microbiome}-- The complex community of trillions of microorganisms, including bacteria, viruses, and fungi, residing in the gastrointestinal tract. The gut microbiome plays a crucial role in digestion, immune function, vitamin synthesis, and protection against pathogens. Its composition is influenced by diet, medications, and lifestyle, and imbalances have been linked to conditions such as obesity, inflammatory bowel disease, and metabolic disorders.

\medterm{Hartmann Procedure}-- Surgical resection of the rectosigmoid colon with creation of an end colostomy and rectal stump. Used for emergency resection of complicated diverticulitis with peritonitis and colorectal cancer with obstruction or perforation. The procedure avoids primary anastomosis in a contaminated field. Reversal with restoration of intestinal continuity is performed electively when the patient recovers, though not all patients are candidates for reversal. Laparoscopic reversal may be technically demanding due to adhesions and fibrosis and carries significant morbidity.

\medterm{Hemochromatosis}-- Iron overload disorder characterized by excessive intestinal iron absorption leading to progressive iron deposition in parenchymal organs. Hereditary hemochromatosis is most commonly caused by homozygous C282Y mutations in the HFE gene affecting approximately one in two hundred individuals of Northern European descent. Clinical manifestations include hepatic cirrhosis, cardiomyopathy, diabetes mellitus, arthropathy, and hypogonadism. Diagnosis is established by elevated transferrin saturation, elevated serum ferritin, and confirmatory genetic testing for HFE mutations.

\medterm{Hemorrhoids}-- Dilated veins in the rectal and anal canal. Internal hemorrhoids are above the dentate line, innervated by visceral nerves, and present with painless bright red bleeding. External hemorrhoids are below the dentate line, innervated by somatic nerves, and may thrombose causing severe pain. Grading: first degree bleed but do not prolapse; second degree prolapse with straining but reduce spontaneously; third degree require manual reduction; fourth degree are irreducible. Managed with dietary fiber, increased fluid intake, stool softeners, and topical treatments. Surgical intervention including haemorrhoidectomy, stapled haemorrhoidopexy, or rubber band ligation is reserved for advanced grades and refractory symptoms.

\medterm{Hepatic Cyst}-- Fluid-filled cavity in the liver most commonly congenital and asymptomatic. Simple hepatic cysts are lined by biliary epithelium and contain clear fluid. May cause pain or compressive symptoms if large or complicated by hemorrhage, infection, or rupture. Simple cysts require no treatment. Polycystic liver disease is an autosomal dominant condition associated with autosomal dominant polycystic kidney disease and may cause significant hepatomegaly and symptoms requiring surgical intervention including fenestration, aspiration sclerotherapy, or surgical resection for symptomatic giant cysts.

\medterm{Hepatic Encephalopathy}-- Neuropsychiatric syndrome occurring in advanced liver disease resulting from accumulation of neurotoxic substances, particularly ammonia. Ammonia crosses the blood-brain barrier and is metabolized by astrocytes into glutamine causing osmotic swelling and impaired neurotransmission. The West Haven grading classifies severity from grade one with mild confusion to grade four with coma. Precipitating factors include gastrointestinal bleeding, infection, constipation, and dehydration. Diagnosis is primarily clinical, supported by elevated serum ammonia, and exclusion of other causes of altered mental status. Management includes lactulose, rifaximin, dietary protein restriction, and treatment of precipitating factors.

\medterm{Hepatic Hemangioma}-- Most common benign liver tumor composed of dilated blood-filled vascular spaces lined by endothelial cells. Usually asymptomatic and found incidentally on imaging. Most are small and require no treatment. Giant hemangiomas exceeding five centimeters may cause pain, compression of adjacent structures, or rarely consumptive coagulopathy with Kasabach-Merritt syndrome. Diagnosed characteristically on imaging showing peripheral nodular enhancement with centripetal fill-in on contrast-enhanced CT or MRI, with a characteristic discontinuous centripetal filling pattern.

\medterm{Hepatic Hydrothorax}-- Large pleural effusion occurring in patients with cirrhosis and portal hypertension in the absence of primary cardiac or pulmonary disease. Most commonly right-sided resulting from transit of ascitic fluid through small defects in the diaphragm into the pleural space driven by the pressure gradient between the abdomen and thorax. Presents with dyspnea, cough, and hypoxemia. Diagnosis requires thoracentesis demonstrating transudative fluid with a high ascitic fluid to pleural fluid protein gradient, along with treatment of the underlying liver disease and portal hypertension.

\medterm{Hepatitis A Virus}-- Picornavirus transmitted through the fecal-oral route most commonly via contaminated food or water. The incubation period is fifteen to fifty days. Clinical presentation ranges from asymptomatic infection in young children to symptomatic acute hepatitis with jaundice, fatigue, nausea, elevated transaminases, dark urine, and pale stools. Unlike hepatitis B and C, hepatitis A does not cause chronic infection. Fulminant hepatitis occurs in approximately zero point one to zero point three percent with overall mortality increasing with age and underlying liver disease.

\medterm{Hepatitis B Virus}-- Hepadnavirus transmitted through percutaneous exposure, sexual contact, and perinatally causing acute and potentially chronic hepatitis. The virus has a partially double-stranded circular DNA genome. Chronic hepatitis B affects approximately two hundred forty million people worldwide, increasing the risk of cirrhosis and hepatocellular carcinoma. Diagnosis is made by serological testing: surface antigen indicates infection, surface antibody indicates immunity, core antibody indicates past or current infection. Vaccination is recommended for all susceptible individuals and is highly effective at preventing infection and its complications.

\medterm{Hepatitis C Virus}-- Flavivirus transmitted primarily through percutaneous exposure with injection drug use being the most common route. Approximately seventy to eighty percent of acute infections progress to chronic hepatitis C, the leading cause of cirrhosis, hepatocellular carcinoma, and liver transplantation in many Western countries. Diagnosis involves anti-hepatitis C antibody screening followed by confirmatory RNA measurement. Treatment with direct-acting antivirals achieves sustained virological response rates exceeding 95 percent, with minimal side effects.

\medterm{Hepatitis D Virus}-- Defective RNA virus requiring co-infection with hepatitis B to replicate, using hepatitis B surface antigen for its envelope. Hepatitis D causes more severe liver disease than hepatitis B alone with accelerated progression to cirrhosis and higher risk of hepatocellular carcinoma. Co-infection with hepatitis D and B during acute hepatitis B is generally self-limited, but superinfection of chronic hepatitis B carriers leads to chronic hepatitis D in seventy to ninety percent of cases. Diagnosis is by serological testing for anti-hepatitis D antibodies and confirmed by HDV RNA detection.

\medterm{Hepatitis E Virus}-- Non-enveloped RNA virus transmitted through the fecal-oral route most commonly via contaminated drinking water. It is the most common cause of acute viral hepatitis worldwide. Clinical presentation is similar to hepatitis A with jaundice and elevated transaminases. The infection is usually self-limited in immunocompetent individuals, but fulminant hepatic failure occurs in approximately one to two percent overall with significantly higher mortality in pregnant women reaching twenty-five percent in the third trimester.

\medterm{Hepatocellular Jaundice}-- Jaundice caused by dysfunction of hepatocytes, impairing the liver's ability to take up, conjugate, or excrete bilirubin. See Jaundice for complete overview. Causes: viral hepatitis (A, B, C, D, E), alcoholic hepatitis, drug-induced liver injury (paracetamol, isoniazid), cirrhosis, autoimmune hepatitis, Wilson disease, haemochromatosis. Laboratory features: mixed elevation of unconjugated and conjugated bilirubin, elevated transaminases (ALT/AST), and impaired synthetic function (low albumin, prolonged INR). Differentiated from obstructive jaundice by normal bile ducts on imaging.

\medterm{Hepatoma}-- See Hepatocellular Carcinoma.

\medterm{Hepatopulmonary Syndrome}-- Triad of liver disease, intrapulmonary vascular dilatation, and impaired oxygenation causing hypoxemia. Presents with platypnea which is dyspnea worsening in the upright position and orthodeoxia which is oxygen desaturation in the upright position. Diagnosed by contrast-enhanced echocardiography showing right-to-left shunt or macroaggregated albumin lung perfusion scan. The condition is a indication for liver transplantation as it may improve or resolve after transplant. There is no effective medical therapy beyond supportive care and oxygen supplementation.

\medterm{Hepatorenal Syndrome}-- Functional renal failure in advanced liver disease characterized by intense renal vasoconstriction in the setting of splanchnic vasodilation. Type 1 hepatorenal syndrome is rapidly progressive with median survival of approximately two weeks without treatment. Type 2 is slowly progressive associated with refractory ascites. Diagnosis requires exclusion of other causes of renal failure. Treatment options include vasoconstrictor therapy with terlipressin and albumin, transjugular intrahepatic portosystemic shunt (TIPS) for refractory cases, and liver transplantation as definitive treatment.

\medterm{Hirschsprung Disease}-- Congenital condition characterized by the absence of enteric ganglion cells in a segment of the distal colon, resulting from failure of neural crest cell migration during embryonic development. The aganglionic segment, which most commonly involves the rectosigmoid colon, remains tonically contracted because it lacks the inhibitory neural input required for relaxation, causing functional bowel obstruction. Affected infants typically present with failure to pass meconium within forty-eight hours of birth, along with abdominal distension, bilious vomiting, and delayed passage of meconium. Diagnosis is confirmed by barium enema showing a transition zone between the dilated proximal ganglionic segment and the narrow distal aganglionic segment, and definitive treatment is surgical resection of the aganglionic segment with pull-through anastomosis.

\medterm{Hirschsprung's Disease}-- A congenital disorder characterized by the absence of parasympathetic ganglion cells (aganglionosis) in the distal colon, resulting in failure of relaxation and functional obstruction. Incidence: 1 in 5,000 live births, male predominance (4:1). The aganglionic segment extends proximally from the internal anal sphincter for a variable distance: short-segment (most common, 80\%, rectosigmoid), long-segment (15\%, extends beyond the sigmoid), total colonic aganglionosis (5\%). Pathophysiology: during embryonic development (weeks 5-12), neural crest cells fail to migrate completely along the hindgut. The aganglionic segment has absent submucosal (Meissner) and myenteric (Auerbach) plexuses, causing tonic contraction and functional obstruction. Mutation in RET proto-oncogene is the most common genetic cause (50\% of familial, 15\% of sporadic cases). Associated with Down syndrome (5-10\%) and other genetic syndromes (Waardenburg, MEN IIA). Clinical presentation: neonates--failure to pass meconium in the first 24-48 hours (99\% of normal infants pass meconium within 24 hours), abdominal distension, bilious vomiting, enterocolitis (fever, explosive diarrhea, sepsis). Older infants/children--chronic constipation, abdominal distension, failure to thrive, enterocolitis. Diagnosis: barium enema (transition zone--narrow aganglionic segment proximal to dilated ganglionic colon), anorectal manometry (absent rectoanal inhibitory reflex), full-thickness rectal biopsy (golden standard--absent ganglion cells, hypertrophied nerve trunks on acetylcholinesterase staining). Treatment: surgical resection of the aganglionic segment with pull-through procedure (Soave, Swenson, Duhamel). Laparoscopic-assisted and transanal endorectal pull-through are preferred. Prognosis: excellent for short-segment disease; long-term issues with constipation, soiling, and enterocolitis in some patients.

\medterm{Ileal Pouch-Anal Anastomosis}-- Surgical creation of a J-pouch from the terminal ileum after proctocolectomy to preserve intestinal continuity for ulcerative colitis. Avoids permanent ileostomy by allowing defecation through the anal route. The procedure involves removal of the entire colon and rectum with construction of a reservoir from distal ileum. Complications include pouchitis, which is the most common complication, cuffitis, anastomotic stricture, and pouch failure. Pouchitis typically responds to antibiotic therapy with metronidazole or ciprofloxacin, and chronic pouchitis may require probiotic therapy, topical steroids, or biologic agents.

\medterm{Inflammatory Bowel Disease}-- Chronic relapsing inflammatory conditions of the gastrointestinal tract encompassing Crohn disease and ulcerative colitis, resulting from a dysregulated immune response in genetically susceptible individuals triggered by environmental factors and alterations in the gut microbiome. The incidence and prevalence of IBD have increased substantially over recent decades, particularly in newly industrialized countries, suggesting important roles for urbanization, dietary changes, and reduced microbial diversity, which collectively drive the chronic inflammatory response.

\medterm{Inflammatory Bowel Disease (IBD)}-- A group of chronic inflammatory conditions affecting the gastrointestinal tract, primarily Crohn's disease and ulcerative colitis. Crohn's disease can involve any part of the digestive tract with transmural inflammation, while ulcerative colitis is confined to the colonic mucosa. Symptoms include abdominal pain, diarrhea, rectal bleeding, and weight loss, managed with anti-inflammatory medications, immunosuppressants, and sometimes surgery.

\medterm{Intestinal Ganglioneuromatosis}-- Benign proliferation of ganglion cells and nerve fibers in the gastrointestinal tract wall. May be diffuse involving the entire gastrointestinal tract or localized as ganglioneuromas. Associated with multiple endocrine neoplasia type 2B where intestinal ganglioneuromatosis may precede the development of medullary thyroid carcinoma. Can cause constipation, diarrhea, abdominal pain, or pseudo-obstruction. Diagnosis is by endoscopic biopsy showing ganglion cells within the lamina propria and muscularis mucosae, often associated with nerve fibre hypertrophy.

\medterm{Intestinal Obstruction}-- A mechanical or functional blockage of intestinal contents. Small bowel obstruction (SBO): most common cause is adhesions (60\%), followed by hernias, tumors, and Crohn disease. Large bowel obstruction (LBO): most common is colorectal cancer, volvulus (sigmoid most common), and diverticulitis. Symptoms: abdominal pain (crampy, colicky), vomiting, abdominal distension, and obstipation. In SBO, vomiting is early and prominent. In LBO, distension is prominent and vomiting is late. Diagnosis: plain abdominal X-ray (air-fluid levels, dilated loops), CT scan (gold standard, identifies etiology and strangulation). Treatment: NPO, NG decompression, IV fluids, and electrolyte correction. Complete SBO often requires surgery (laparoscopy or laparotomy with adhesiolysis). LBO: resection, Hartmann procedure, or colonic stenting as bridge to surgery.

\medterm{Intestinal Pseudo-Obstruction}-- Acute or chronic dilation of the bowel without mechanical obstruction due to neuromuscular dysfunction. Acute colonic pseudo-obstruction also called Ogilvie syndrome occurs in hospitalized or post-surgical patients. Chronic intestinal pseudo-obstruction may be idiopathic or secondary to systemic diseases including scleroderma, diabetes, and Parkinson disease. Managed with decompression, prokinetic agents including metoclopramide and erythromycin, and treatment of underlying conditions. Neostigmine is the first-line pharmacological treatment for acute colonic pseudo-obstruction, with resolution rates exceeding 80 percent, but requires cardiac monitoring due to risk of bradycardia.

\medterm{Intra-Abdominal Abscess}-- Localized collection of pus within the abdomen commonly in the pelvis, between bowel loops, or beneath the diaphragm. Results from intra-abdominal infection or surgery. Presents with fever, abdominal pain, and leukocytosis. Diagnosed by CT with contrast showing rim-enhancing fluid collection. Treated with percutaneous CT-guided drainage and broad-spectrum antibiotics. Surgery is reserved for multiloculated abscesses not amenable to percutaneous drainage or failure of percutaneous management.

\medterm{Intra-Abdominal Hypertension}-- Sustained elevation of intra-abdominal pressure above twelve millimeters of mercury measured through the urinary bladder. Graded as grade one twelve to fifteen, grade two sixteen to twenty, grade three twenty-one to twenty-five, and grade four greater than twenty-five millimeters of mercury. Causes include fluid resuscitation, abdominal surgery, ileus, ascites, and abdominal hemorrhage. Intra-abdominal hypertension can progress to abdominal compartment syndrome with organ dysfunction. Monitoring of intra-abdominal pressure is essential in critically ill and post-surgical patients to guide interventions and prevent progression to abdominal compartment syndrome.

\medterm{Intussusception}-- Condition in which a segment of intestine telescopes into the adjacent distal segment causing bowel obstruction and compromise of the mesenteric blood supply. It is the most common cause of intestinal obstruction in infants and children between six months and three years of age. Classic presentation includes episodic colicky abdominal pain, vomiting, a sausage-shaped abdominal mass, and current jelly stool. Diagnosis is confirmed by ultrasound demonstrating a target sign or donut sign. In stable patients with non-ischemic intussusception, air or contrast enema reduction is both diagnostic and therapeutic, with success rates exceeding 80 percent in children.

\medterm{Irritable Bowel Syndrome}-- Functional gastrointestinal disorder characterized by recurrent abdominal pain associated with altered bowel habits without identifiable structural abnormalities. The Rome IV criteria define IBS as recurrent abdominal pain occurring at least one day per week associated with changes in defecation or stool form. Subtypes include IBS with constipation, IBS with diarrhea, and IBS with mixed bowel habits. Pathophysiology involves visceral hypersensitivity, altered gut motility, gut-brain axis dysfunction, and psychosocial factors. Management includes dietary modifications, stress reduction, fibre supplementation, antispasmodics, and neuromodulators.

\medterm{Irritable Bowel Syndrome (IBS)}-- A chronic functional gastrointestinal disorder characterized by recurrent abdominal pain and altered bowel habits (diarrhea, constipation, or mixed) without identifiable structural pathology. Pathophysiology involves dysregulation of the gut--brain axis, visceral hypersensitivity, altered motility, and imbalances in the gut microbiome. Diagnosis is clinical based on Rome IV criteria: recurrent abdominal pain at least 1 day/week for 3 months, associated with defecation or change in stool frequency/form. Management includes dietary modifications (low-FODMAP diet), fiber supplementation, antispasmodics, probiotics, and neuromodulators (tricyclic antidepressants, SSRIs) for refractory symptoms.

\medterm{Ischemic Colitis}-- Colonic mucosal ischemia resulting from reduced blood flow through the marginal artery and end-arteries, most commonly occurring at watershed areas between vascular territories, particularly the splenic flexure. Unlike acute mesenteric ischemia affecting the small bowel, ischemic colitis most often occurs with non-occlusive hypoperfusion from cardiac failure or dehydration. The clinical presentation includes sudden onset crampy left-sided abdominal pain followed by bloody diarrhea. Most cases are self-limiting and resolve with supportive care and bowel rest, with full recovery expected within one to two weeks without specific intervention.

\medterm{Jaundice}-- Yellowish discoloration of the skin, sclerae, and mucous membranes caused by hyperbilirubinemia (serum bilirubin greater than 2-3 mg/dL), resulting from an imbalance between bilirubin production and clearance. Bilirubin metabolism: heme from senescent red blood cells is converted by heme oxygenase to biliverdin, then to unconjugated (indirect) bilirubin, which is transported to the liver bound to albumin. Unconjugated bilirubin is taken up by hepatocytes and conjugated with glucuronic acid by uridine diphosphate glucuronosyltransferase (UGT) in the hepatocyte. Conjugated (direct) bilirubin is water-soluble and excreted into bile. Hyperbilirubinaemia is classified as predominantly unconjugated (pre-hepatic or hepatic) or conjugated (hepatic or post-hepatic/obstructive).

\medterm{Jaundice (Icterus)}-- Yellowish discoloration of the skin, sclerae, and mucous membranes caused by hyperbilirubinemia (serum bilirubin greater than 2-3 mg/dL), resulting from an imbalance between bilirubin production and clearance. See \hyperlink{Jaundice}{Jaundice}.

\medterm{Kayser-Fleischer Rings}-- Golden-brown copper deposition rings in Descemet membrane of the cornea representing a pathognomonic finding of Wilson disease. These rings are visible on slit lamp examination and present in virtually all patients with neurological manifestations of Wilson disease though they may be absent in isolated hepatic disease. The rings typically begin at the superior and inferior poles and progress circumferentially. They may gradually diminish after initiation of chelation therapy. Detection requires slit-lamp examination by an experienced ophthalmologist.

\medterm{Lactose}-- See biochemistry.

\medterm{Lactose intolerance}-- A common digestive disorder caused by lactase deficiency leading to malabsorption of lactose, the primary disaccharide in milk and dairy products. The lactase enzyme is produced by enterocytes of the small intestinal brush border and hydrolyzes lactose into glucose and galactose for absorption. Three types exist: primary lactase deficiency (most common, involves the genetically programmed decline in lactase activity after weaning, affecting 70-75 percent of the world's adult population, with highest prevalence in individuals of East Asian, African, and Middle Eastern descent where lactase persistence is less common.

\medterm{Large Bowel Obstruction}-- Mechanical blockage of the colon most commonly caused by colorectal cancer, volvulus, or diverticular stricture. Clinical presentation includes progressive abdominal distension, constipation, and obstipation. The cecum is at highest risk of perforation with distal obstruction. Diagnosis is by CT scan identifying the site and cause. Management depends on cause: sigmoid volvulus may be managed by endoscopic detorsion, cecal volvulus requires emergent surgery, and colorectal cancer obstruction may require a diverting colostomy or colonic stent as a bridge to elective resection.

\medterm{Liver Abscess}-- Focal collection of pus within the liver parenchyma. Pyogenic abscesses are most commonly from biliary tract infection or hematogenous spread from portal or systemic infection. Amoebic abscesses are caused by Entamoeba histolytica predominantly in tropical regions. Presents with fever, right upper quadrant pain, and hepatomegaly. Diagnosed by CT showing rim-enhancing fluid collection. Treatment involves percutaneous CT-guided drainage and broad-spectrum antibiotics for pyogenic abscesses. Metronidazole is the treatment of choice for amoebic liver abscesses, with aspiration reserved for large abscesses at risk of rupture.

\medterm{Liver Biopsy}-- Percutaneous, transjugular, or laparoscopic sampling of liver tissue for histological evaluation. Indicated for staging fibrosis, diagnosing specific liver diseases, and assessing treatment response. The procedure is guided by ultrasound for accuracy. Complications include bleeding, pain, and pneumothorax. Contraindications include uncorrected coagulopathy and thrombocytopenia. Transjugular biopsy is performed when percutaneous biopsy is contraindicated due to coagulopathy or ascites. Histologic evaluation remains the gold standard for assessing fibrosis stage and grade of inflammatory activity in chronic liver disease.

\medterm{Liver Cirrhosis}-- Advanced hepatic fibrosis with regenerative nodules surrounded by dense fibrous septa representing the end result of chronic liver injury. Common etiologies include chronic viral hepatitis, alcoholic liver disease, and nonalcoholic steatohepatitis. The pathogenesis involves activation of hepatic stellate cells producing excessive extracellular matrix. Cirrhosis progresses through compensated and decompensated phases. Decompensation manifests with ascites, variceal hemorrhage, encephalopathy, and ultimately hepatic failure.

\medterm{Liver Diseases}-- A broad category of pathological conditions affecting the structure and function of the liver. See individual entries: Hepatitis, Liver Cirrhosis, Fatty Liver, Hepatocellular Carcinoma, etc.

\medterm{Liver Health}-- Essential for metabolism, detoxification, protein synthesis, bile production, and immune function. Common liver disorders include nonalcoholic fatty liver disease (NAFLD, now termed metabolic dysfunction--associated steatotic liver disease or MASLD) affecting ~30\% of adults, alcoholic liver disease (steatosis, hepatitis, cirrhosis), viral hepatitis (HAV, HBV, HCV), drug-induced liver injury (acetaminophen toxicity most common), and autoimmune hepatitis. MASLD progresses from steatosis to steatohepatitis (NASH/MASH) with fibrosis and cirrhosis, driven by insulin resistance, obesity, and metabolic syndrome. Screening: liver enzymes (ALT, AST), viral serologies, abdominal ultrasound, FibroScan for fibrosis, and liver biopsy when indicated. Management: weight loss $\ge$7\% improves steatohepatitis in MASLD, alcohol cessation, antiviral therapy for HBV/HCV (direct-acting antivirals cure >95\% of HCV), and vaccination for HAV/HBV.

\medterm{Liver Transplantation}-- Surgical replacement of a diseased liver with a healthy donor organ indicated for end-stage liver disease, acute liver failure, and select hepatocellular carcinomas within Milan criteria. Requires lifelong immunosuppression with calcineurin inhibitors, antimetabolites, and corticosteroids. Living donor liver transplantation uses a portion of a living donor's liver which regenerates in both donor and recipient. Complications include rejection, infection, vascular complications, biliary complications (anastomotic strictures, bile leaks), and hepatic artery thrombosis.

\medterm{Low FODMAP Diet}-- Dietary intervention for irritable bowel syndrome that restricts fermentable oligosaccharides, disaccharides, monosaccharides, and polyols, which are short-chain carbohydrates poorly absorbed in the small intestine and rapidly fermented by colonic bacteria producing gas and drawing water into the intestinal lumen. The approach involves an elimination phase of two to six weeks, a reintroduction phase identifying personal triggers, and a long-term personalization phase. Studies demonstrate symptom improvement in 50-80 percent of patients with irritable bowel syndrome.

\medterm{Lymphocytic Colitis}-- Chronic inflammatory bowel disease characterized by increased intraepithelial lymphocytes exceeding twenty per one hundred surface epithelial cells in the colonic mucosa without significant subepithelial collagen band thickening. Shares clinical features with collagenous colitis including chronic watery non-bloody diarrhea primarily in middle-aged women. The pathogenesis may involve an immune response to luminal antigens. Diagnosis requires colonoscopy with random colonic biopsies showing increased intraepithelial lymphocytes on histology.

\medterm{Malabsorption}-- Defective absorption of nutrients from the small intestine. Causes: celiac disease, pancreatic insufficiency, bacterial overgrowth, Crohn disease, short bowel syndrome, and infections. Presents with diarrhea, steatorrhea, weight loss, and deficiency of specific nutrients (iron, B12, vitamin D, calcium). Diagnosis: stool studies, D-xylose test, breath tests, small bowel biopsy. Treatment: address underlying cause and replace deficient nutrients.

\medterm{Meconium Ileus}-- Intestinal obstruction in neonates caused by inspissated meconium most commonly associated with cystic fibrosis. Presents with abdominal distension, bilious vomiting, and failure to pass meconium within forty-eight hours. Diagnosis is suggested by abdominal radiography showing a ground-glass appearance and may be confirmed by CT showing calcified meconium. Treatment includes water-soluble contrast enema which can be both diagnostic and therapeutic by loosening the inspissated meconium. Surgical intervention is required when non-operative management fails, with laparotomy and enterotomy for meconium removal and possible bowel resection for complicated cases.

\medterm{MELD Score}-- Model for End-Stage Liver Disease predicting three-month mortality in advanced liver disease used for liver transplant allocation. The score is calculated using serum bilirubin, creatinine, and international normalized ratio generating a continuous score from approximately six to forty points. Higher scores indicate more severe disease and higher mortality risk. The updated MELD score includes serum sodium. MELD exception points are granted for conditions including hepatocellular carcinoma within the Milan criteria (single tumor up to 5 cm, or up to 3 tumors each up to 3 cm, without extrahepatic spread or major vascular invasion).

\medterm{Microscopic Colitis}-- Umbrella term encompassing both collagenous colitis and lymphocytic colitis which share the clinical presentation of chronic watery non-bloody diarrhea and the pathophysiological feature of colonic mucosal inflammation. Endoscopic appearance is typically normal necessitating random biopsies for diagnosis. More common in women and associated with autoimmune diseases, celiac disease, and certain medications including proton pump inhibitors, selective serotonin reuptake inhibitors, and nonsteroidal anti-inflammatory drugs (NSAIDs), which can cause mucosal injury through inhibition of cyclooxygenase and prostaglandin synthesis.

\medterm{Mirizzi Syndrome}-- Condition where a gallstone impacted in the cystic duct or Hartmann pouch extrinsic compression of the common hepatic duct causing obstructive jaundice. Classified into four types based on the degree of erosion and fistula formation. Type one is extrinsic compression without fistula. Types two through four involve cholecystobiliary fistula of increasing size. Presents with obstructive jaundice, cholangitis, and abdominal pain. Diagnosis is by MRCP or endoscopic ultrasound demonstrating the impacted gallstone on the bile duct, and cholecystectomy with fistula repair is the definitive management.

\medterm{MR Enterography}-- MRI of the small bowel after oral contrast administration. Superior soft tissue contrast for evaluating Crohn disease activity, fistulas, and abscesses without ionizing radiation. Demonstrates mural thickening, edema, enhancement patterns, and extraintestinal complications. Diffusion-weighted imaging helps assess disease activity. Preferred imaging modality for young patients requiring serial follow-up due to absence of radiation. The technique is limited by longer examination time, availability, and susceptibility to motion artefact.

\medterm{Neuroendocrine Tumors}-- Heterogeneous group of neoplasms arising from neuroendocrine cells throughout the gastrointestinal tract and pancreas. May be functional producing hormones or non-functional. Graded by Ki-67 index and mitotic count into well-differentiated low-grade, intermediate-grade, and poorly differentiated high-grade tumors. Functional tumors include gastrinomas, insulinomas, glucagonomas, VIPomas, and carcinoid tumors secreting serotonin. Diagnosis is by hormone levels, imaging with CT, MRI, or octreotide scintigraphy (OctreoScan), and tissue biopsy for histological grading. Management includes somatostatin analogues for symptom control, peptide receptor radionuclide therapy (PRRT), and surgical resection for localised disease.

\medterm{Non-Alcoholic Fatty Liver Disease}-- Hepatic steatosis in the absence of significant alcohol consumption, strongly associated with metabolic syndrome, obesity, and type 2 diabetes. NAFLD includes simple steatosis and non-alcoholic steatohepatitis (NASH) with inflammation and fibrosis. NASH can progress to cirrhosis and hepatocellular carcinoma. Diagnosis: abdominal ultrasound, elevated liver enzymes, transient elastography, and liver biopsy for staging. Management: weight loss, exercise, vitamin E, and pioglitazone for NASH.

\medterm{Nonalcoholic Steatohepatitis}-- Progressive form of NAFLD characterized by hepatic steatosis with lobular inflammation and hepatocyte ballooning with significant risk of progression to fibrosis, cirrhosis, and hepatocellular carcinoma. NASH affects approximately three to five percent of the general population. The histological features include macrovesicular steatosis, inflammation, hepatocyte ballooning, and perisinusoidal fibrosis. Progressive fibrosis is the most important predictor of liver-related mortality. Management involves lifestyle modification with weight loss, dietary changes, exercise, vitamin E therapy in non-diabetic patients with biopsy-proven NASH, and consideration of obeticholic acid or GLP-1 receptor agonists in selected patients.

\medterm{Obstructive Jaundice}-- Jaundice caused by obstruction of bile flow from the liver to the duodenum, resulting in accumulation of conjugated bilirubin in the blood. See Jaundice for complete overview. Causes: (1) Extrahepatic obstruction: choledocholithiasis (common bile duct stones, most common), pancreatic head carcinoma (painless jaundice, Courvoisier law), cholangiocarcinoma, ampullary tumor, primary sclerosing cholangitis, biliary stricture (post-surgical, post-ERCP), Mirizzi syndrome. (2) Intrahepatic obstruction: primary biliary cholangitis (now PBC), drug-induced cholestasis, infiltrative diseases (sarcoidosis, amyloidosis), sepsis, total parenteral nutrition. Laboratory features: elevated conjugated (direct) bilirubin, raised ALP and GGT (cholestatic enzymes), transaminases may be normal or mildly elevated. Imaging: ultrasound (dilated bile ducts, stone, mass), MRCP, ERCP (therapeutic). Complications: cholangitis, secondary biliary cirrhosis, fat-soluble vitamin deficiency (A, D, E, K), prpruritus.

\medterm{Pancreatic Cancer}-- Aggressive malignancy with pancreatic ductal adenocarcinoma accounting for approximately eighty-five percent of cases. It is the fourth leading cause of cancer-related death with a five-year survival of approximately ten percent. Risk factors include smoking, chronic pancreatitis, diabetes, and genetic syndromes. The most common presentation is painless obstructive jaundice from head of pancreas tumors. The Courvoisier sign suggests malignant biliary obstruction. Diagnosis involves CT, endoscopic ultrasound with fine-needle aspiration for tissue diagnosis, and tumor marker CA 19-9. Staging uses CT and endoscopic ultrasound.

\medterm{Paracentesis}-- Percutaneous needle aspiration of ascitic fluid from the peritoneal cavity. Used for diagnostic analysis to determine the cause of ascites and for therapeutic removal of large volumes for symptom relief. The serum-ascites albumin gradient is calculated to differentiate portal hypertensive from non-portal hypertensive causes. A gradient of one point one grams per deciliter or greater indicates portal hypertension. Diagnostic paracentesis with cell count, culture, and cytology should be performed before initiating treatment in all patients with new-onset ascites.

\medterm{Peak Flow Meter}-- See Peak Expiratory Flow Rate.

\medterm{Peritoneal Carcinomatosis}-- Diffuse metastatic seeding of the peritoneal cavity from gastrointestinal or ovarian malignancies. Presents with ascites, bowel obstruction, and cachexia. Diagnosed by CT and paracentesis showing malignant cells. Prognosis is generally poor. Management includes systemic chemotherapy, palliative care, and in selected cases, cytoreductive surgery with hyperthermic intraperitoneal chemotherapy. The peritoneal cancer index assesses disease distribution for surgical planning.

\medterm{Peritoneal Dialysis Catheter}-- Flexible tube surgically placed into the peritoneal cavity for peritoneal dialysis. The Tenckhoff catheter is the most commonly used type with one or two cuffs that promote tissue ingrowth and anchor the catheter. Placement can be done surgically, laparoscopically, or percutaneously under radiological guidance. Complications include peritonitis, catheter-related infection, tunnel infection, outflow failure, and hernia formation. Proper catheter placement technique and exit site care are essential for preventing infectious complications and ensuring catheter longevity.

\medterm{Pilonidal Disease}-- A chronic inflammatory condition of the sacrococcygeal region involving hair follicles. Pilonidal sinus: epithelialized tract containing hair and debris. Pilonidal abscess: acute infection of the sinus. Most common in young, hirsute males, exacerbated by prolonged sitting, obesity, and local trauma. Presents with pain, swelling, and purulent drainage in the natal cleft. Treatment: incision and drainage for acute abscess; definitive surgical treatment (excision with secondary healing, flap closure, or pit-picking) for chronic disease. Recurrence is common without good hygiene and hair removal.

\medterm{Pneumatosis Cystoides Intestinalis}-- Gas-filled cysts within the bowel wall that may be primary or secondary to bowel ischemia, infection, chronic lung disease, or medications. Often an incidental finding on imaging. Primary pneumatosis is benign and asymptomatic. Secondary pneumatosis may indicate serious underlying disease. Managed by treating the underlying condition. Surgery is reserved for complications including pneumoperitoneum, obstruction, or hemorrhage.

\medterm{Pneumatosis Intestinalis}-- Presence of gas within the bowel wall appearing as linear or cystic radiolucencies on imaging. Primary pneumatosis is benign and idiopathic accounting for approximately fifteen percent of cases. Secondary pneumatosis is associated with intestinal ischemia, necrotizing enterocolitis in neonates, obstructive pulmonary disease, and connective tissue diseases. The cystic form appears as clusters of gas-filled cysts in the submucosa or subserosa. The linear form suggests intramural gas tracking along the mesenteric vessels, which may indicate bowel necrosis and requires urgent surgical intervention.

\medterm{Polyp Surveillance Colonoscopy}-- Interval-based endoscopic follow-up of patients with previously identified colorectal polyps. Intervals are determined by polyp number, size, and histological features. One to two small adenomas require surveillance at seven to ten years. Three to four adenomas at three to five years. Five to ten adenomas or adenomas with villous features or high-grade dysplasia at three years. More than ten adenomas at one year. Sessile serrated polyps greater than ten millimeters require three to five year surveillance intervals.

\medterm{Porcelain Gallbladder}-- Calcification of the gallbladder wall visible on imaging representing chronic inflammatory condition associated with cholelithiasis and chronic cholecystitis. Historically considered a strong risk factor for gallbladder carcinoma, recent studies suggest the cancer risk may be lower than previously estimated. Most patients are asymptomatic and the condition is discovered incidentally. When cholecystectomy is performed, the gallbladder should be carefully examined for occult malignancy. Management is cholecystectomy, preferably by the laparoscopic approach.

\medterm{Postcholecystectomy Syndrome}-- Persistence or recurrence of symptoms after cholecystectomy. Causes include retained common bile duct stones, sphincter of Oddi dysfunction, bile salt malabsorption, irritable bowel syndrome, and other gastrointestinal conditions. Evaluation includes liver function tests, MRCP, and endoscopic ultrasound. Treatment depends on the identified cause. Bile salt malabsorption responds to bile acid sequestrants. Sphincter of Oddi dysfunction may require endoscopic sphincterotomy.

\medterm{Prebiotics and Probiotics}-- Prebiotics are nondigestible food ingredients (fiber and oligosaccharides—inulin, fructooligosaccharides, galactooligosaccharides) that selectively stimulate the growth and activity of beneficial gut bacteria (Bifidobacterium, Lactobacillus), promoting host health. Probiotics are live microorganisms that confer a health benefit when administered in adequate amounts—commonly \emph{Lactobacillus} species, \emph{Bifidobacterium}, \emph{Saccharomyces boulardii} (yeast), and \emph{Streptococcus thermophilus}. Evidence supports probiotics for prevention of antibiotic-associated diarrhea (~50\% risk reduction), reduction of \emph{C. difficile} infection recurrence (alongside standard antibiotics), management of irritable bowel syndrome (bloating, pain), treatment of acute infectious gastroenteritis (reduces duration by ~1 day), and prevention of necrotizing enterocolitis in preterm infants. Synbiotics combine prebiotics and probiotics. The microbiome modulates immune function, metabolism, and gut--brain signaling. Quality control varies widely; products should specify genus, species, strain, and colony-forming units (CFU) per dose. Fermented foods (yogurt, kefir, kimchi, sauerkraut) naturally provide beneficial microbes.

\medterm{Primary Biliary Cholangitis}-- Chronic autoimmune cholestatic liver disease characterized by progressive destruction of small intrahepatic bile ducts. The disease predominantly affects middle-aged women. Anti-mitochondrial antibodies are found in approximately ninety-five percent of patients. The hallmark histological finding is florid duct lesion with granulomatous destruction of interlobular bile ducts. Clinical presentation includes fatigue and prpruritus. Diagnosis is established by elevated alkaline phosphatase, positive anti-mitochondrial antibodies, elevated serum IgM, and characteristic histology on liver biopsy. First-line treatment is ursodeoxycholic acid, which improves liver biochemistry and delays disease progression.

\medterm{Primary Sclerosing Cholangitis}-- Chronic cholestatic liver disease characterized by progressive fibrosis and strictures of intrahepatic and extrahepatic bile ducts. The disease has a strong association with ulcerative colitis occurring in seventy to eighty percent of patients. Diagnosis is established by MRCP showing multifocal strictures and dilatations giving a beaded appearance. Patients have a significantly increased risk of cholangiocarcinoma developing in ten to fifteen percent. There is no proven effective medical therapy, and liver transplantation is the definitive treatment for end-stage disease.

\medterm{Proctitis}-- Inflammation of the rectal mucosa causing tenesmus, rectal pain, and bloody discharge. Causes include inflammatory bowel disease, sexually transmitted infections including gonorrhea, chlamydia, and herpes, radiation therapy, and ischemia. Diagnosed by proctoscopy with biopsy. Treatment is directed at the underlying etiology. Empiric treatment for sexually transmitted infections should be initiated while awaiting culture results. Inflammatory proctitis responds to topical mesalamine or corticosteroid enemas.

\medterm{Pruritus Ani}-- Persistent itching around the anus with multiple potential causes including hemorrhoids, fissures, infections such as pinworms and candidiasis, dermatologic conditions including psoriasis and lichen sclerosus, and idiopathic causes. Managed by addressing the underlying cause, strict hygiene measures with gentle cleansing after bowel movements, avoiding scratching, and topical treatments including hydrocortisone cream and barrier preparations. Cotton underwear and avoidance of irritating foods such as coffee, spicy foods, and citrus fruits should also be advised.

\medterm{Pseudomembranous Colitis}-- Colonic inflammation caused by C. difficile infection showing yellow-white plaques on colonoscopy. The plaques consist of fibrin, inflammatory cells, and necrotic epithelial debris overlying mucosal ulceration. Presents with profuse watery diarrhea and abdominal pain. Diagnosis is confirmed by endoscopy with biopsy showing characteristic pseudomembranes. Treated with oral vancomycin or fidaxomicin. Severe cases may require intravenous metronidazole and surgical intervention for fulminant colitis unresponsive to medical therapy.

\medterm{Psoas Abscess}-- Collection of pus in the iliopsoas muscle compartment. Primary psoas abscess results from hematogenous spread of Staphylococcus aureus typically in immunocompromised patients. Secondary psoas abscess results from contiguous spread from adjacent structures including Crohn disease, diverticulitis, appendicitis, and vertebral osteomyelitis. Presents with fever, flank or abdominal pain, and characteristic hip flexion posture. Diagnosis is by CT of the abdomen and pelvis showing the psoas collection. Treatment is CT-guided percutaneous drainage and intravenous antibiotics directed at the causative organism.

\medterm{Push Enteroscopy}-- Endoscopic technique using a longer endoscope to visualize and intervene in the proximal small bowel beyond the reach of standard upper endoscopy. Unlike capsule endoscopy, push enteroscopy allows direct visualization with the ability to obtain biopsies and perform hemostasis. The procedure is performed with the endoscope advanced through the mouth into the proximal and mid-jejunum. The primary indication is evaluation of obscure gastrointestinal bleeding identified on capsule endoscopy. The diagnostic yield of push enteroscopy is highest when performed within 24-48 hours of active bleeding, and it enables targeted therapeutic intervention.

\medterm{Ranson Criteria}-- Clinical scoring system for acute pancreatitis severity incorporating eleven parameters at two time points. Admission parameters include age, white blood cell count, glucose, LDH, and AST. Forty-eight hour parameters include hematocrit drop, BUN increase, calcium, oxygen tension, base deficit, and fluid sequestration. A score of three or more indicates severe pancreatitis with mortality increasing with score. The Ranson criteria require forty-eight hours to complete, making them less useful for early severity assessment, and the Bedside Index for Severity in Acute Pancreatitis (BISAP) score is increasingly preferred for early risk stratification.

\medterm{Rectal Prolapse}-- Protrusion of rectal tissue through the anal opening. Complete prolapse involves all layers of the rectal wall and presents as a concentric cylindrical mass with visible mucosal folds. Partial prolapse involves only the mucosa. Associated with chronic straining, pelvic floor dysfunction, and neurological conditions. Treatment is surgical with rectopexy being the preferred approach. Perineal approaches including Delorme procedure and Altemeier procedure are options for elderly or high-risk patients.

\medterm{Rectovaginal Fistula}-- Abnormal connection between the rectum and vagina causing passage of stool or gas through the vagina. Causes include obstetric injury particularly third and fourth degree perineal tears, Crohn disease, radiation therapy, and surgical complications. Presents with vaginal discharge of stool, flatus, or recurrent vaginal infections. Classified by location as low involving the anal sphincter or high above the sphincter. Treated surgically with approaches including transanal advancement flap, transperineal repair, or with gracilis muscle flap interposition for complex recurrent fistulae.

\medterm{Reynolds Pentad}-- Extension of the Charcot triad to include hypotension and altered mental status representing suppurative cholangitis, a severe life-threatening form of ascending cholangitis. Patients present with high fever, rigors, right upper quadrant pain, jaundice, hypotension, and confusion. Management involves aggressive fluid resuscitation, vasopressor support, broad-spectrum antibiotics, and emergent biliary decompression by ERCP. Without prompt treatment, suppurative cholangitis carries mortality exceeding 50 percent without prompt intervention.

\medterm{Rovsing Sign}-- Physical exam sign of acute appendicitis: palpation of the left lower quadrant elicits pain in the right lower quadrant due to displacement of gas stretching the inflamed appendix. \textit{See Appendicitis.}

\medterm{Salmonellosis}-- Infection with Salmonella bacteria (non-typhoidal), typically acquired from contaminated food (eggs, poultry, meat). Presents with acute gastroenteritis: diarrhea, fever, and abdominal cramps within 6-72 hours of ingestion. Self-limited in immunocompetent patients. Bacteremia can occur, especially in immunocompromised. Treatment: supportive care; antibiotics reserved for severe disease, bacteremia, or immunocompromised hosts.

\medterm{Serrated Pathway}-- Alternative molecular pathway to colorectal carcinogenesis accounting for fifteen to thirty percent of cases. It involves serrated polyps as precursors including sessile serrated lesions with dysplasia. Molecular features include CpG island methylator phenotype, BRAF mutations, and microsatellite instability from MLH1 silencing. Sessile serrated lesions tend to be flat, sessile, mucus-covered, and located in the proximal colon making them difficult to detect. Recognition has led to updated colonoscopic screening and surveillance guidelines that emphasize careful inspection of the proximal colon.

\medterm{Shigellosis}-- Acute infectious colitis caused by Shigella species (S. sonnei most common in developed countries). Presents with fever, abdominal cramps, tenesmus, and frequent bloody, mucoid stools. Highly contagious with low infectious dose. Diagnosis: stool culture. Treatment: azithromycin, ciprofloxacin, or ceftriaxone (resistance patterns vary). Complications: reactive arthritis, hemolytic uremic syndrome (S. dysenteriae type 1).

\medterm{Short Bowel Syndrome}-- Clinical condition resulting from surgical resection or functional loss of a significant portion of the small intestine leading to inadequate absorption of nutrients. Clinical manifestations depend on the length and location of the remaining bowel, the presence or absence of the ileocecal valve, and the adaptability of the remaining intestine. Patients typically experience profuse diarrhea, dehydration, electrolyte imbalances, and progressive malnutrition. Those losing the terminal ileum lack vitamin B12 absorption, requiring lifelong parenteral supplementation, while those with an intact colon can absorb short-chain fatty acids and salvage some calories and fluid.

\medterm{Spontaneous Bacterial Peritonitis}-- Infection of ascitic fluid without an identifiable intra-abdominal source representing a serious complication of cirrhosis with ascites. The pathogenesis involves impaired immune function, bacterial translocation from the intestinal lumen, and seeding of ascitic fluid which provides a favorable medium for bacterial growth. The most common causative organism is Escherichia coli. Patients present with fever, abdominal pain, and altered mental status, though symptoms may be absent in thirty percent of cases, necessitating a high index of suspicion and a low threshold for diagnostic paracentesis.

\medterm{Tenesmus}-- A sensation of incomplete defecation, characterized by a constant feeling of needing to pass stool accompanied by straining, pain, and cramping, despite little or no stool being passed. Tenesmus is a symptom, not a diagnosis. Causes: inflammatory bowel disease (ulcerative colitis—most common, Crohn disease), infectious colitis (Shigella, Salmonella, Campylobacter, Clostridioides difficile, amoebic dysentery), radiation proctitis (post-pelvic radiotherapy), colorectal cancer (especially rectal cancer—should be excluded), anorectal conditions (hemorrhoids, anal fissure, solitary rectal ulcer syndrome, proctalgia fugax), and irritable bowel syndrome. Pathophysiology: inflammation, infiltration, or irritation of the rectal mucosa and nerve endings (S2--S4, pelvic splanchnic nerves) triggers the urge to defecate. Assessment: proctoscopy/sigmoidoscopy or colonoscopy with biopsy, stool culture, and imaging (CT, MRI) if tumor suspected. Treatment: directed at the underlying cause—mesalazine, steroids, and biologic therapy for IBD; antibiotics for infectious colitis; surgery for malignancy; topical anaesthetics, biofeedback, and amitriptyline for functional disorders.

\medterm{Toxic Megacolon}-- Acute life-threatening dilation of the colon associated with systemic toxicity, representing one of the most serious complications of inflammatory bowel disease, particularly ulcerative colitis. It can also occur in infectious colitis caused by Clostridioides difficile and cytomegalovirus. The pathophysiology involves transmural inflammation extending beyond the mucosa, causing destruction of the myenteric plexus and smooth muscle dysfunction leading to loss of colonic tone. The diagnosis requires radiologic evidence of colonic dilation exceeding 6 cm in diameter, with loss of haustrations, accompanied by clinical signs of toxicity including fever, tachycardia, leukocytosis, and systemic inflammatory response.

\medterm{Tropical Sprue}-- Chronic diarrhea and malabsorption syndrome of uncertain etiology occurring predominantly in tropical and subtropical regions. The condition is characterized by villous atrophy of the small intestinal mucosa leading to malabsorption of nutrients including folate, iron, and vitamin B12. The exact pathogenesis remains unclear but may involve an infectious etiology, altered gut microbiota, or nutritional deficiencies. The clinical presentation includes chronic diarrhea, steatorrhea, weight loss, and anemia. Treatment includes a prolonged course of tetracycline or doxycycline with folic acid supplementation, and the condition responds well to antibiotic therapy if diagnosed early.

\medterm{Typhlitis}-- Inflammation of the cecum in immunocompromised patients particularly those with neutropenia from chemotherapy. Also known as neutropenic enterocolitis. Presents with right lower quadrant pain, fever, and diarrhea. Diagnosis is by CT showing cecal wall thickening. Managed with bowel rest, broad-spectrum antibiotics covering gram-negative and anaerobic organisms, and resolution of neutropenia. Surgery is reserved for complications including perforation and hemorrhage.

\medterm{Typhoid Fever}-- A systemic infection caused by Salmonella typhi, acquired through contaminated food or water. Incubation 7-14 days. Presents with prolonged fever, headache, abdominal pain, rose spots, and relative bradycardia. Diagnosis: blood culture (first week), stool culture (second week), Widal test. Treatment: azithromycin, ceftriaxone, or fluoroquinolones. Complications: intestinal perforation, hemorrhage. Prevention: vaccination and water sanitation.

\medterm{Ulcerative Colitis}-- Chronic inflammatory disease limited to the colon and rectum, characterized by continuous mucosal and submucosal inflammation that invariably begins in the rectum and extends proximally in an uninterrupted pattern. The distribution follows a predictable anatomical pattern classified as proctitis when confined to the rectum, left-sided colitis extending to the splenic flexure, or extensive colitis extending beyond the splenic flexure. Unlike Crohn disease, the inflammation in ulcerative colitis is confined to the mucosa and submucosa, without transmural involvement.

\medterm{Volvulus}-- Twisting of a loop of intestine around its mesenteric attachment point causing luminal obstruction and compromise of the mesenteric blood supply which can rapidly progress to bowel ischemia. The two most common sites are the sigmoid colon accounting for approximately seventy-five percent of colonic volvulus in elderly patients with chronic constipation and redundant sigmoid, and the cecum accounting for approximately twenty-five percent in younger patients with a mobile cecum. Sigmoid volvulus presents with abdominal distension, pain, and obstipation, and is initially managed with endoscopic detorsion using a sigmoidoscope, followed by elective sigmoid resection to prevent recurrence.

\medterm{Whipple Disease}-- A rare systemic bacterial infection caused by Tropheryma whipplei, affecting the small intestine and multiple organ systems. Presents with arthralgias (prodromal, years before GI symptoms), weight loss, chronic diarrhea, steatorrhea, and abdominal pain. Extraintestinal: endocarditis, neurologic symptoms (dementia, ophthalmoplegia, myoclonus), and lymphadenopathy. Diagnosis: small bowel biopsy showing PAS-positive macrophages containing rod-shaped bacteria, PCR of tissue or CSF. Treatment: prolonged antibiotic therapy (ceftriaxone IV induction, followed by TMP-SMX or doxycycline + hydroxychloroquine for 1-2 years). Untreated is fatal.

\medterm{Whipple Triad}-- Diagnostic criteria for Insulinoma: (1) Symptoms of hypoglycemia (diaphoresis, confusion), (2) Documented low plasma glucose ($<55\text{ mg/dL}$), and (3) Prompt relief of symptoms following glucose administration. \textit{See Insulinoma.}

\medterm{Zieve's Syndrome}-- A triad of haemolytic anemia, hyperlipidaemia, and jaundice occurring in the setting of chronic alcohol use disorder and alcoholic liver disease (acute alcoholic steatohepatitis or alcoholic cirrhosis). First described by Leslie Zieve in 1958 in patients with alcoholic liver disease presenting with transient haemolytic anemia (Coombs-negative, likely due to oxidative damage to erythrocyte membranes from acetaldehyde and lipid peroxidation products), markedly elevated serum triglycerides (hyperlipidaemia, often severe, >1000 mg/dL, due to alcohol-induced increase in hepatic VLDL secretion and reduced lipoprotein lipase activity), and jaundice (from alcoholic hepatitis, intrahepatic cholestasis, and haemolysis). The syndrome is typically self-limited, resolving with alcohol abstinence over 2-4 weeks, but recurrence is common if alcohol consumption resumes. Management: supportive care, alcohol cessation, and management of hypertriglyceridaemia (dietary fat restriction, fibrates) to prevent acute pancreatitis.

\medterm{Chronic Liver Disease}-- Persistent hepatic injury and structural change lasting at least 6 months, caused by viral, metabolic, alcohol-related, autoimmune, cholestatic, vascular, or drug-related disease. Progressive fibrosis can lead to cirrhosis, portal hypertension, liver failure, and hepatocellular carcinoma.

\medterm{Decompensated Cirrhosis}-- Cirrhosis complicated by clinically significant liver dysfunction or portal hypertension, commonly manifested by ascites, variceal hemorrhage, hepatic encephalopathy, or jaundice. It carries substantially higher mortality than compensated cirrhosis and warrants assessment for transplant and complication prevention.

\medterm{Hepatic Synthetic Dysfunction}-- Impaired hepatic production of proteins and other molecules, including albumin and coagulation factors. It may cause hypoalbuminemia, edema, prolonged prothrombin time, and increased INR, although INR can also be altered by vitamin K deficiency, anticoagulants, and other conditions.

\medterm{Portal Vein Thrombosis}-- Partial or complete thrombosis of the portal vein or its branches. It may result from cirrhosis, malignancy, inflammation, thrombophilia, or abdominal surgery and can cause portal hypertension, intestinal ischemia, or variceal formation.


\medterm{Biliary Colic}-- Episodic right-upper-quadrant or epigastric pain caused by transient obstruction of the cystic duct, usually by a gallstone. Pain is typically steady, may follow meals, and resolves when the obstruction passes; fever and persistent inflammatory findings suggest a complication.

\medterm{Acute Cholecystitis}-- Acute inflammation of the gallbladder, most often caused by persistent cystic-duct obstruction by a gallstone. It commonly causes sustained right-upper-quadrant pain, fever, nausea, and a positive Murphy sign; ultrasound and timely antimicrobial and surgical assessment are standard.

\medterm{Pancreatic Exocrine Insufficiency}-- Inadequate secretion or delivery of pancreatic digestive enzymes and bicarbonate, causing maldigestion and malabsorption. Chronic pancreatitis, cystic fibrosis, pancreatic resection, and advanced pancreatic disease are important causes; steatorrhea and weight loss may respond to pancreatic enzyme replacement.

\medterm{Pancreatic Pseudocyst}-- A mature, encapsulated collection of pancreatic fluid and inflammatory debris without a true epithelial lining, usually developing after pancreatitis or pancreatic trauma. Symptomatic, infected, enlarging, or complicated pseudocysts may require endoscopic, percutaneous, or surgical drainage.

\medterm{Secretory Diarrhea}-- Diarrhea caused by active secretion or reduced absorption of electrolytes and water by the intestine. It persists despite fasting and may result from infections, medications, bile-acid disorders, endocrine tumors, or inflammatory disease.

\medterm{Osmotic Diarrhea}-- Diarrhea caused by poorly absorbed solutes retaining water in the intestinal lumen. It typically improves with fasting and may result from carbohydrate malabsorption, laxatives, or other osmotically active substances.

\medterm{Acute-on-Chronic Liver Failure (ACLF)}-- Acute deterioration of pre-existing chronic liver disease associated with organ failure and high short-term mortality. It is commonly precipitated by infection, alcohol-related hepatitis, bleeding, or other acute insults and requires urgent multidisciplinary management.


\medterm{Gastrointestinal Perforation}-- A full-thickness defect in the gastrointestinal wall allowing luminal contents to enter the peritoneal or adjacent space. It can cause peritonitis, sepsis, and shock and usually requires urgent imaging, antimicrobial therapy, resuscitation, and source control.
\medterm{Achalasia}-- See \hyperlink{Esophageal Achalasia}{Esophageal Achalasia}.
\medterm{Acute esophageal necrosis}-- Ischemic necrosis of the esophageal mucosa, most commonly caused by low-flow states such as shock, sepsis, or thromboembolism, presenting with hematemesis and chest pain. Endoscopy shows circumferential black necrotic mucosa, typically in the distal esophagus. Management is supportive with acid suppression, and mortality is high due to underlying comorbidities.
\medterm{Bezoar}-- See \hyperlink{Gastric Bezoar}{Gastric Bezoar}.
\medterm{Bilateral peritonsillar abscesses}-- Suppurative infection in the peritonsillar space on both sides, a rare complication of acute tonsillitis presenting with severe sore throat, dysphagia, trismus, and muffled hot potato voice. Bilateral involvement risks acute upper airway obstruction and requires urgent needle aspiration or incision and drainage. Associated with fever, odynophagia, and potential spread to deep neck spaces.
\medterm{Bladder diverticula}-- Outpouchings of the bladder mucosa through the muscular wall, often acquired from chronic bladder outlet obstruction or congenital in origin. Acquired diverticula are associated with benign prostatic hyperplasia and may cause urinary stasis, recurrent infections, and stone formation. Congenital diverticula, such as Hutch diverticula near the ureterovesical junction, may predispose to vesicoureteral reflux.
\medterm{Bowel infarction}-- Ischemic necrosis of bowel tissue from occlusive or non-occlusive mesenteric ischemia, presenting with severe abdominal pain out of proportion to examination, bloody stools, and peritoneal signs. Causes include mesenteric arterial embolism or thrombosis, venous thrombosis, and low-flow states. Mortality is high, and emergent surgical resection of necrotic bowel is required.
\medterm{Bronchoesophageal fistula}-- Abnormal communication between the bronchial tree and the esophagus, most commonly acquired from malignancy, mediastinal infection, or tuberculous lymphadenitis. Congenital forms are associated with esophageal atresia. Presents with recurrent aspiration pneumonia, coughing and choking with oral intake, and recurrent lower lobe infiltrates. Diagnosis is by CT chest with contrast and confirmed by upper endoscopy.
\medterm{Cecal volvulus}-- Axial twisting of the cecum and ascending colon around its mesentery, accounting for approximately one-third of colonic volvulus cases. It presents with acute or subacute large bowel obstruction, abdominal pain, and distension. Diagnosis is by CT scan showing the characteristic whirl sign and a massively dilated cecum. Treatment requires emergent surgery with cecectomy or right hemicolectomy due to high risk of gangrene and perforation.
\medterm{Colocolonic intussusception}-- telescoping of a segment of colon into an adjacent segment, less common in adults than the ileocolic type and usually associated with a pathological lead point such as a polyp, tumor, or lipoma. Adults typically present with intermittent abdominal pain, bloating, and change in bowel habits. Diagnosis is by CT showing the target sign or sausage-shaped mass, and treatment involves surgical resection.
\medterm{Colon cancer}-- Malignant neoplasm arising from the colonic epithelium, the third most common cancer worldwide and a leading cause of cancer-related mortality. Most arise from adenomatous polyps through the adenoma-carcinoma sequence over ten to fifteen years. Risk factors include older age, inflammatory bowel disease, hereditary polyposis syndromes, and dietary factors. Diagnosis is by colonoscopy with biopsy, and staging uses TNM classification.
\medterm{Duodenal perforation}-- Full-thickness breach of the duodenal wall, most commonly caused by penetrating peptic ulcer disease, trauma, or iatrogenic injury during endoscopy. Presents with sudden severe epigastric pain, board-like abdominal rigidity, and peritonitis. Free air under the diaphragm on upright chest radiograph is diagnostic. Emergent surgical repair with omental patch closure is required.
\medterm{Encapsulating peritoneal sclerosis}-- Rare condition characterized by a thick fibro-collagenous membrane encasing the small intestine, leading to recurrent bowel obstruction. Most commonly associated with long-term peritoneal dialysis but also seen after abdominal surgery, radiation, and beta-blocker use. Presents with abdominal pain, distension, nausea, and intermittent obstruction. Diagnosis is by CT showing cocoon-like encasement of bowel loops.
\medterm{Esophagogastric bypass}-- Surgical creation of a conduit between the esophagus and stomach to restore gastrointestinal continuity, most commonly performed as part of esophagectomy or to palliate unresectable esophageal malignancy. Techniques include gastric pull-up with esophagogastrostomy and jejunal interposition. Complications include anastomotic leak, stricture, and conduit necrosis.
\medterm{Gallbladder volvulus}-- Axial torsion of the gallbladder around its mesentery and cystic duct, most commonly affecting elderly women, leading to acute gallbladder ischemia. Presents with sudden onset severe right upper quadrant pain, nausea, vomiting, and signs of peritoneal irritation. Diagnosis is by ultrasound showing a distended, non-functioning gallbladder, and treatment requires emergent cholecystectomy.
\medterm{Gastric cancer}-- See \hyperlink{Stomach Cancer}{Stomach Cancer}.
\medterm{Gastric pneumatosis}-- Presence of gas within the gastric wall, also called gastric emphysema, which may be benign (gastric emphysema) or life-threatening (emphysematous gastritis). Benign forms result from barotrauma, chronic obstructive pulmonary disease, or mechanical causes. Emphysematous gastritis is caused by gas-forming organisms and presents with severe abdominal pain, sepsis, and gastric wall necrosis requiring emergent intervention.
\medterm{Gastrogastric intussusception}-- Rare form of intussusception where one segment of the stomach telescopes into an adjacent gastric segment, typically associated with a lead point such as a gastric polyp, leiomyoma, or carcinoma. Presents with abdominal pain, nausea, vomiting, and occasionally upper gastrointestinal bleeding. Diagnosis is by CT scan or upper endoscopy, and treatment involves surgical or endoscopic reduction with resection of the lead point.
\medterm{Gastrointestinal amyloidosis}-- Deposition of misfolded amyloid fibrils in the gastrointestinal tract, either as part of systemic amyloidosis (AL, AA, or ATTR types) or rarely as primary localized disease. Involvement can occur anywhere from the oral cavity to the rectum, causing dysmotility, malabsorption, ischemic ulceration, protein-losing enteropathy, or gastrointestinal bleeding. Diagnosis requires tissue biopsy with Congo red staining showing apple-green birefringence under polarized light.
\medterm{Gastrospenic fistulization}-- Abnormal communication between the stomach and spleen, most commonly arising from gastric ulcer erosion into the splenic hilum or as a complication of splenic abscess. Presents with abdominal pain, gastrointestinal hemorrhage, and sepsis. Diagnosis is by CT with oral contrast showing extravasation, and management requires emergent surgical repair with possible splenectomy.
\medterm{Infantile hepatic hemangiomas}-- Benign vascular tumors of the liver occurring in infants, typically presenting within the first six months of life with hepatomegaly, high-output cardiac failure, consumptive hypothyroidism, or cutaneous hemangiomas. Large tumors may cause abdominal compartment syndrome and thrombocytopenia from Kasabach-Merritt phenomenon. Diagnosis is by imaging with characteristic peripheral enhancement on MRI, and management includes propranolol as first-line medical therapy.
\medterm{Intestinal lipomatosis}-- Benign overgrowth of adipose tissue within the intestinal wall, most commonly affecting the ileum and colon. Often asymptomatic and discovered incidentally, but large lipomas can cause intussusception, gastrointestinal bleeding, or bowel obstruction. Diagnosis is by CT or endoscopy showing well-defined fatty masses, and symptomatic lesions are treated by endoscopic or surgical resection.
\medterm{Intraabdominal hernia}-- Herniation of abdominal viscera through a congenital or acquired defect or abnormal mesenteric aperture within the peritoneal cavity, distinct from ventral or inguinal hernias. Types include paraduodenal, foramen of Winslow, transmesocolic, and transomental hernias. Often present with intermittent or acute bowel obstruction, and diagnosis is by CT scan. Surgical reduction and closure of the defect is required to prevent strangulation.
\medterm{Lung herniation}-- Protrusion of lung tissue through a defect in the chest wall, most commonly occurring after thoracic trauma, thoracotomy, or in association with rib fractures and chronic obstructive pulmonary disease. Presents with a reducible chest wall mass that increases with coughing or straining and may cause pain and respiratory compromise. Diagnosis is by CT chest, and symptomatic hernias require surgical repair with mesh closure.
\medterm{Malakoplakia}-- Rare chronic inflammatory condition characterized by accumulation of macrophages containing calcified inclusion bodies called Michaelis-Gutmann bodies, most commonly affecting the urinary bladder but also occurring in the gastrointestinal tract. It is associated with impaired macrophage bactericidal function and occurs predominantly in immunocompromised patients. Diagnosis is by biopsy showing characteristic Michaelis-Gutmann bodies on histology.
\medterm{Meckel's diverticulum}-- See \hyperlink{Meckel Diverticulum}{Meckel Diverticulum}.
\medterm{Metastasis of colorectal cancer}-- Spread of colorectal malignancy to distant organs via lymphatic, hematogenous, or peritoneal routes, most commonly involving the liver, lungs, and peritoneum. The liver is the first site of metastasis due to portal venous drainage, and hepatic metastases may be resectable with curative intent. Staging with CT, MRI, and PET guides management, which may include systemic chemotherapy, targeted biologic therapy, hepatic resection, or ablation.
\medterm{Metastatic colon cancer}-- Advanced colonic malignancy that has spread beyond the bowel to distant sites, most commonly the liver, lungs, and peritoneum, via lymphatic and hematogenous dissemination. Staging uses CT of chest, abdomen, and pelvis, and PET-CT may identify occult metastases. Management involves palliative systemic chemotherapy, targeted therapy with anti-VEGF or anti-EGFR agents, and possible surgical resection of oligometastatic disease.
\medterm{Neurogenic megacolon}-- Massive dilation of the colon due to loss of autonomic innervation, most commonly from Chagas disease, Hirschsprung disease, or diabetic autonomic neuropathy. It presents with chronic constipation, abdominal distension, fecal impaction, and risk of toxic megacolon. Diagnosis is by abdominal radiography showing colonic dilation and colonic manometry, and management may include colonoscopic decompression or surgical resection.
\medterm{Nutcracker esophagus}-- Esophageal motility disorder characterized by high-amplitude peristaltic contractions exceeding 180 mmHg in the distal esophagus on manometry, with preserved peristaltic coordination. Presents with intermittent chest pain and dysphagia to solids and liquids. Diagnosis is by high-resolution manometry using Chicago classification criteria, and management includes calcium channel blockers, nitrates, and anxiolytics.
\medterm{Obturator hernia}-- Herniation of abdominal contents through the obturator canal, most commonly occurring in thin elderly women with a history of multiple pregnancies. It is a rare cause of small bowel obstruction and may present with the Howship-Romberg sign of medial thigh pain from obturator nerve compression. Diagnosis is difficult clinically and is usually made by CT scan, and surgical repair is required to prevent bowel strangulation.
\medterm{Paraesophageal hernia}-- See \hyperlink{Hiatal Hernia}{Hiatal Hernia}.
\medterm{Perforation of a viscus}-- Full-thickness breach of the wall of a hollow abdominal organ, most commonly the stomach, duodenum, or colon, resulting in leakage of luminal contents into the peritoneal cavity. Causes include peptic ulcer disease, diverticulitis, appendicitis, inflammatory bowel disease, ischemia, and trauma. Presents with acute severe abdominal pain, peritonitis, and free air on imaging. Emergent surgical repair is required.
\medterm{Peripancreatic malignancy}-- Neoplastic process involving tissues surrounding the pancreas, including primary pancreatic adenocarcinoma with peripancreatic extension and metastatic disease from adjacent or distant organs. Often presents at an advanced stage with abdominal pain, weight loss, obstructive jaundice, and new-onset diabetes. Diagnosis is by CT or MRI with biopsy when indicated, and management depends on resectability.
\medterm{Pneumatosis cystoides coli}-- Condition characterized by multiple gas-filled cysts within the wall of the colon, which may be primary (idiopathic) or secondary to intestinal mucosal disruption from conditions such as chronic obstructive pulmonary disease, ischemic bowel disease, or intestinal obstruction. Often discovered incidentally on imaging or colonoscopy. Management is directed at the underlying cause, and the condition may resolve spontaneously.
\medterm{Pneumoperitoneum}-- Free air within the peritoneal cavity, most commonly caused by perforation of a hollow viscus such as the stomach, duodenum, or colon. Iatrogenic causes include recent laparoscopy, endoscopy with perforation, and peritoneal dialysis. Presents with acute abdominal pain, rigidity, and signs of peritonitis. Diagnosis is by upright chest radiograph showing air under the diaphragm or by CT scan, and emergent surgical repair of the perforation is required.
\medterm{Pseudomelanosis duodeni}-- Benign condition characterized by dark pigment deposits in the duodenal mucosa, visible as slate-gray or black patches on endoscopy, caused by accumulation of iron-containing pigment in macrophages within the lamina propria. Associated with chronic kidney disease, gastrointestinal bleeding, and certain medications such as ferrous sulfate and hydrochlorothiazide. It is a histological finding of no clinical significance.
\medterm{Pylephlebitis}-- Septic thrombophlebitis of the portal vein or its tributaries, most commonly arising from intra-abdominal infections such as appendicitis, diverticulitis, cholangitis, or hepatic abscess. Presents with high fever, rigors, right upper quadrant pain, jaundice, and sepsis. Diagnosis is by CT with contrast showing portal vein thrombosis, and treatment involves broad-spectrum antibiotics and management of the underlying source.
\medterm{Retroperitoneal hemorrhage}-- Bleeding into the retroperitoneal space, which may result from aortic rupture, anticoagulation therapy, retroperitoneal hematoma following femoral catheterization, trauma, or retroperitoneal neoplasms. Presents with flank pain, abdominal distension, hypotension, and anemia. Diagnosis is by CT scan showing blood in the retroperitoneal compartment, and management includes resuscitation, reversal of coagulopathy, and treatment of the underlying cause.
\medterm{Sigmoid volvulus}-- Axial twisting of the sigmoid colon around its mesentery, the most common cause of colonic volvulus, typically occurring in elderly, institutionalized, or neuropsychiatric patients with chronic constipation and a redundant sigmoid. Presents with acute large bowel obstruction, abdominal distension, and absent bowel sounds. Diagnosis is by CT showing the whirl sign, and treatment involves endoscopic detorsion followed by elective sigmoid resection to prevent recurrence.
\medterm{Small-bowel volvulus}-- Twisting of the small intestine around its mesentery, causing mechanical obstruction and potentially compromising mesenteric blood flow. It may be primary (idiopathic, more common in developing countries) or secondary to adhesions, malrotation, hernias, or intraluminal tumors. Presents with colicky abdominal pain, vomiting, distension, and signs of strangulation. Diagnosis is by CT, and emergent surgery is required to reduce the volvulus and resect necrotic bowel.
\medterm{Tracheal agenesis with bronchoesophageal fistula}-- Congenital malformation characterized by absent or severely hypoplastic trachea with an abnormal communication between the bronchus and esophagus, a rare and typically fatal condition presenting in neonates with respiratory distress, cyanosis, and inability to intubate. Associated anomalies include esophageal atresia, cardiac defects, and vertebral anomalies. Diagnosis is by CT and bronchoscopy, and management requires emergent surgical repair.
\medterm{Trichobezoar}-- See \hyperlink{Gastric Bezoar}{Gastric Bezoar}.
\medterm{Zenker's diverticulum}-- See \hyperlink{Esophageal Diverticulum}{Esophageal Diverticulum}.
\medterm{Villous atrophy}-- Flattening and loss of the normal small-intestinal villi, reducing absorptive surface area and occurring in conditions such as coeliac disease.
\medterm{Haemorrhoids}-- Enlarged symptomatic vascular cushions of the anal canal that may cause painless bright-red rectal bleeding, prolapse, itching, or discomfort.

\medterm{Cullen's Sign}-- Bluish or violaceous periumbilical ecchymosis caused by blood tracking through the fascial planes to the umbilicus. It is a sign of intra-abdominal or retroperitoneal haemorrhage, classically severe haemorrhagic pancreatitis, and is associated with serious underlying disease.

\medterm{Grey Turner's Sign}-- Bluish or violaceous ecchymosis of the flanks caused by retroperitoneal bleeding tracking through the abdominal wall. It is classically associated with severe haemorrhagic pancreatitis but may occur with other retroperitoneal haemorrhage.

\medterm{Abdominal Distention}-- Visible or subjective enlargement, tightness, or fullness of the abdomen caused by gas, fluid, stool, pregnancy, organ enlargement, mass, or bowel obstruction. Acute distention with pain, vomiting, obstipation, fever, haemodynamic instability, or peritoneal signs requires urgent evaluation.

\medterm{Abdominal Mass}-- A palpable or imaging-detected abnormal collection, enlargement, or lesion within the abdomen. Causes include distended bowel or bladder, organ enlargement, hernia, inflammatory collection, benign tumour, and malignancy; assessment depends on location, mobility, tenderness, systemic features, and appropriate imaging.

\medterm{Acute Hepatitis}-- Acute inflammation and injury of hepatocytes, usually developing over days to months and characterized by elevated aminotransferases, with or without jaundice, nausea, right-upper-quadrant discomfort, or constitutional symptoms. Causes include viral infection, autoimmune disease, ischaemia, toxins, and metabolic or other inflammatory injury.

\medterm{Biliary Cirrhosis}-- See \hyperlink{Primary Biliary Cholangitis}{Primary Biliary Cholangitis}.

\medterm{Diaphragmatic Hernia}-- Protrusion of abdominal contents through a defect or opening in the diaphragm into the thoracic cavity. It may be congenital or acquired after trauma and can cause respiratory compromise, bowel obstruction, or strangulation; diagnosis is based on imaging and management depends on symptoms, contents, and the risk of organ compromise.

\medterm{Enterocolitis}-- Inflammation of both the small intestine and colon, caused by infection, immune-mediated disease, reduced blood flow, toxins, or other injury. It may produce abdominal pain, diarrhoea, fever, vomiting, bleeding, dehydration, or systemic illness; diagnosis and management depend on the cause and severity.

\medterm{Fecal Impaction}-- A large, firm mass of stool retained in the rectum or colon that cannot be evacuated normally. It may cause abdominal or rectal pain, distention, overflow diarrhoea, urinary symptoms, nausea, or bowel obstruction and is associated with constipation, immobility, dehydration, neurological disease, and medications.

\medterm{Flatulence}-- Passage of intestinal gas through the anus. It reflects swallowed air and bacterial fermentation of undigested carbohydrates and may be increased by dietary factors, constipation, malabsorption, altered gut motility, or functional gastrointestinal disorders; persistent symptoms require assessment for associated disease.

\medterm{Gastroesophageal Reflux}-- Backward flow of gastric contents into the oesophagus, caused by transient or persistent relaxation of the lower oesophageal sphincter or other mechanical factors. It may cause heartburn, regurgitation, chest or upper-abdominal discomfort, cough, hoarseness, or dental erosion; frequent symptoms may represent gastroesophageal reflux disease.

\medterm{Gastrointestinal Obstruction}-- Partial or complete blockage of the passage of intestinal contents through the gastrointestinal tract. It may be mechanical or functional and can cause colicky pain, vomiting, distention, constipation, obstipation, dehydration, bowel ischaemia, perforation, or sepsis.

\medterm{Hepatosplenomegaly}-- Enlargement of both the liver and spleen. It may result from infection, portal hypertension, storage disease, malignancy, haemolysis, or systemic inflammatory and infiltrative disorders; evaluation depends on associated symptoms and laboratory and imaging findings.

\medterm{Ileal Atresia}-- Congenital complete interruption or absence of a segment of the ileal lumen. It causes neonatal intestinal obstruction with abdominal distension, bilious vomiting, and failure to pass meconium and generally requires surgical management.

\medterm{Ileitis}-- Inflammation of the ileum, the distal segment of the small intestine. Causes include Crohn disease, infection, ischemia, medication-related injury, and other inflammatory disorders; symptoms may include right-lower-quadrant pain, diarrhea, fever, or bleeding.

\medterm{Ileocecal Valve}-- A sphincter-like fold at the junction of the terminal ileum and caecum. It regulates passage of intestinal contents into the colon and helps limit reflux of colonic contents into the small intestine.

\medterm{Ileus}-- Failure of coordinated intestinal propulsion without a mechanical blockage, usually from postoperative bowel inhibition, inflammation, metabolic disturbance, or medication effect. It causes distension, nausea, vomiting, and reduced passage of flatus or stool.

\medterm{Jejunal Atresia}-- Congenital complete interruption or absence of a segment of the jejunal lumen. It presents in a newborn with intestinal obstruction, commonly causing bilious vomiting, abdominal distension, and failure to pass meconium, and generally requires surgery.

\medterm{Linitis Plastica}-- A diffuse infiltrative form of gastric carcinoma in which tumour cells spread through the stomach wall, producing marked thickening and rigidity with loss of normal distensibility. It may cause early satiety, weight loss, nausea, vomiting, or vague upper-abdominal symptoms.

\medterm{Radiation Enteritis}-- Inflammation and injury of the small or large intestine caused by therapeutic ionizing radiation. It may be acute with diarrhea, cramping, and nausea or delayed with strictures, malabsorption, bleeding, fistula, or obstruction.

\medterm{Wandering Spleen}-- Abnormal mobility of the spleen caused by laxity or absence of its supporting ligaments. It may be asymptomatic or cause abdominal pain, splenic enlargement, torsion, infarction, or acute abdomen.

\medterm{Helicobacter Pylori Infection}-- Infection of the stomach by H. pylori, a spiral-shaped bacterium associated with chronic gastritis, peptic ulcer disease, gastric adenocarcinoma, and gastric MALT lymphoma. It is diagnosed with noninvasive or endoscopic testing and treated with combination eradication therapy.

\medterm{Hiccups}-- Recurrent involuntary spasms of the diaphragm followed by sudden closure of the glottis, producing a characteristic sound. Persistent or intractable hiccups may reflect irritation of the diaphragm or phrenic and vagal pathways, metabolic disease, medication exposure, or central nervous system disease.
